\chapter{Teoria pierścieni: od arytmetyki do geometrii równań}
\label{ch:pierscienie}
\section{Po co nam pierścienie?}

Przed badaniem zbiorów opisanych równaniami wielomianowymi potrzebujemy języka,
który pozwoli rozumieć same wielomiany i relacje między nimi. Nie chodzi tylko
 o wykonywanie rachunków. Równania $f=0$ oraz $g=0$ mogą opisywać ten sam zbiór,
różne zbiory, albo różne \emph{struktury algebraiczne} ukryte pod tym samym
zbiorem punktów. Te różnice najlepiej badać za pomocą pierścieni i ideałów.

\begin{przyklad}[Trzy rodzaje arytmetyki]
Rozważmy wyrażenie $x^2-1=(x-1)(x+1)$.
\begin{enumerate}[label=\arabic*)]
\item W $\R$ z równania $x^2=1$ wynika $x=1$ albo $x=-1$.
\item W $\Z/8\Z$ równanie $x^2=1$ ma aż cztery rozwiązania:
$\overline1,\overline3,\overline5,\overline7$.
Rzeczywiście, jeśli $x=2m+1$, to $x^2=4m(m+1)+1$, a ponieważ $m(m+1)$
jest parzyste, mamy $x^2\equiv1\pmod8$.
Jeżeli natomiast $x$ jest parzyste, to $x^2$ jest parzyste,
a więc nie może być przystające do $1$ modulo $8$.
Wszystkie cztery klasy nieparzyste i tylko one są zatem rozwiązaniami.
\item W pierścieniu $\R[t]/(t^2)$ element $t$ jest niezerowy, lecz
$t^2=0$. Nie wolno zatem wnioskować z $ab=0$, że $a=0$ lub $b=0$.
\end{enumerate}
Te różnice nie wynikają z odmiennego zapisu symboli, lecz z własności
struktury, w której wykonujemy działania.
\end{przyklad}

\begin{przyklad}[Równania i funkcje]
Wielomian $y-x^2\in\R[x,y]$ opisuje parabolę. Gdy na tej paraboli
pracujemy z wielomianami, nie potrzebujemy osobno pamiętać wszystkich
wielomianów przyjmujących na niej te same wartości: dwa z nich różnią się
o wielokrotność $y-x^2$ dokładnie wtedy, gdy wyznaczają na \emph{rzeczywistych}
punktach paraboli tę samą funkcję wielomianową. Wykorzystujemy tu, że
niezerowy wielomian jednej zmiennej nad $\R$ nie może mieć nieskończenie
wielu pierwiastków; nad ciałem skończonym analogiczne twierdzenie o funkcjach
nie jest prawdziwe. Wyrazimy to później izomorfizmem
$\R[x,y]/(y-x^2)\cong\R[x]$ i udowodnimy go, zamiast jedynie odwoływać się
do rysunku.
\end{przyklad}

W tym rozdziale budujemy pojęcia w następującej kolejności:
\emph{pierścień $\to$ elementy szczególne $\to$ homomorfizm $\to$
ideał $\to$ pierścień ilorazowy $\to$ ideał pierwszy i maksymalny $\to$
pierścienie wielomianów}. Na końcu połączymy je z pierwszymi przykładami
geometrycznymi. Nie zakładamy znajomości teorii pierścieni; potrzebna jest
jedynie elementarna algebra i pojęcia grupy oraz odwzorowania.

\section{Definicja pierścienia i jej konsekwencje}

\begin{definicja}[Pierścień]
\label{def:pierscien}
Pierścieniem nazywamy zbiór $A$ z działaniami
$+\colon A\times A\to A$ i $\cdot\colon A\times A\to A$ oraz wyróżnionymi
elementami $0,1\in A$, które spełniają warunki:
\begin{enumerate}[label=(P\arabic*)]
\item $(A,+)$ jest grupą abelową: dodawanie jest łączne i przemienne,
$0$ jest elementem neutralnym, a każdy $a\in A$ ma element przeciwny $-a$;
\item mnożenie jest łączne: $(ab)c=a(bc)$ dla wszystkich $a,b,c\in A$;
\item $1$ jest obustronnym elementem neutralnym mnożenia:
$1a=a1=a$ dla każdego $a\in A$;
\item zachodzą obie rozdzielności:
$a(b+c)=ab+ac$ oraz $(a+b)c=ac+bc$ dla wszystkich $a,b,c\in A$.
\end{enumerate}
Pierścień jest \emph{przemienny}, gdy ponadto $ab=ba$ dla dowolnych $a,b\in A$.
\end{definicja}

Pozwalamy na $1=0$; wtedy, jak zaraz sprawdzimy, pierścień ma tylko jeden
element i nazywa się \emph{pierścieniem zerowym}. Jest to wygodne, ponieważ
iloraz $A/A$ będzie pierścieniem. \textbf{Jeżeli nie powiedziano inaczej,
od tego miejsca wszystkie pierścienie są przemienne i mają jedynkę.}
Gdy mowa o dziedzinie lub ciele, wymagamy dodatkowo $1\ne0$.

\begin{uwaga}[Dlaczego rozdzielamy aksjomaty?]
W definicji nie ma założenia, że wolno dzielić przez element różny od zera.
Nie zakładamy nawet, że z $ab=ac$ i $a\ne0$ wynika $b=c$.
Te własności odróżnią później ciała i dziedziny od ogólnych pierścieni.
Również przemienność mnożenia nie wynika z pozostałych aksjomatów:
na przykład macierze kwadratowe tworzą zwykle pierścień nieprzemienny.
\end{uwaga}

\begin{lemat}[Podstawowe reguły rachunku]
\label{lem:reguly}
Niech $A$ będzie dowolnym pierścieniem, niekoniecznie przemiennym.
Dla wszystkich $a,b\in A$ zachodzą równości
$0a=a0=0$, $(-a)b=-(ab)=a(-b)$ oraz $(-a)(-b)=ab$.
Jeżeli $1=0$, to $A=\{0\}$.
\end{lemat}
\begin{proof}
Z rozdzielności mamy $0a=(0+0)a=0a+0a$.
Dodając do obu stron element przeciwny do $0a$ w grupie $(A,+)$,
otrzymujemy $0=0a$. Analogicznie $a0=0$.
Ponieważ $(-a)b+ab=((-a)+a)b=0b=0$, element $(-a)b$ jest przeciwny
do $ab$, czyli $(-a)b=-(ab)$. Tak samo $a(-b)=-(ab)$.
Wreszcie $(-a)(-b)=-\bigl(a(-b)\bigr)=-(-(ab))=ab$.
Jeżeli $1=0$, to dla każdego $a\in A$ mamy $a=1a=0a=0$.
\end{proof}

\begin{definicja}[Podpierścień]
Niech $A$ będzie pierścieniem. Podzbiór $B\subseteq A$ jest
\emph{podpierścieniem}, jeżeli $0,1\in B$, dla wszystkich $a,b\in B$
zachodzi $a-b\in B$ i $ab\in B$. Działania na $B$ są ograniczeniami
działań w $A$.
\end{definicja}

Warunek $1\in B$ jest częścią naszej konwencji. Dlatego $2\Z$ nie jest
podpierścieniem $\Z$ (chociaż jest zamknięty na dodawanie, odejmowanie
i mnożenie). Później okaże się, że $2\Z$ jest \emph{ideałem}.

\begin{przyklad}[Łańcuch podpierścieni]
Mamy $\Z\subseteq\Q\subseteq\R\subseteq\C$.
Każdy z tych zbiorów zawiera tę samą jedynkę $1$ i jest zamknięty
na odpowiednie działania. Podpierścienie mogą się znacznie różnić:
w $\Z$ liczba $2$ nie ma odwrotności względem mnożenia,
a w $\Q$ taką odwrotnością jest $1/2$.
\end{przyklad}

\section{Najważniejsze przykłady pierścieni}

\subsection{Liczby całkowite i arytmetyka modulo \texorpdfstring{$n$}{n}}

W $\Z$ dodajemy i mnożymy w zwykły sposób. Aby otrzymać nowy przykład,
ustalamy $n\ge2$ i mówimy, że $a,b\in\Z$ są przystające modulo $n$,
jeżeli $n\mid(a-b)$. Zapisujemy to $a\equiv b\pmod n$.
Relacja ta jest relacją równoważności: zwrotność wynika z $n\mid0$,
symetria z $n\mid(a-b)\Rightarrow n\mid(b-a)$,
a przechodniość z faktu, że suma dwóch wielokrotności $n$
również jest wielokrotnością $n$.

Klasę liczby $a$ oznaczamy $\overline a$; zbiór klas to $\Z/n\Z$.
Chcemy zdefiniować $\overline a+\overline b=\overline{a+b}$ i
$\overline a\,\overline b=\overline{ab}$. Trzeba sprawdzić,
że wynik nie zależy od wybranych reprezentantów.

\begin{stwierdzenie}
\label{st:modulo}
Dla każdego $n\ge2$ powyższe wzory definiują działania na $\Z/n\Z$,
a z tymi działaniami $\Z/n\Z$ jest pierścieniem przemiennym.
\end{stwierdzenie}
\begin{proof}
Załóżmy, że $a'=a+nr$ oraz $b'=b+ns$ dla pewnych $r,s\in\Z$.
Wtedy $(a'+b')-(a+b)=n(r+s)$, zatem $a'+b'\equiv a+b\pmod n$.
Ponadto
$a'b'-ab=(a+nr)(b+ns)-ab=n(as+br+nrs)$,
więc $a'b'\equiv ab\pmod n$. Działania są zatem dobrze określone.
Wszystkie aksjomaty pierścienia przechodzą z $\Z$ na klasy:
przykładowo $\overline a(\overline b+\overline c)
=\overline{a(b+c)}=\overline{ab+ac}
=\overline a\,\overline b+\overline a\,\overline c$.
Pozostałe aksjomaty sprawdza się w ten sam sposób, używając
odpowiednich własności liczb całkowitych.
\end{proof}

\begin{center}
\begin{tikzpicture}[x=0.79cm,y=1cm,>=Stealth]
\draw[->,gray!75] (-6.7,0)--(7.1,0) node[right] {$\Z$};
\foreach \i in {-6,...,6} {\draw (\i,0.10)--(\i,-0.10);}
\foreach \i in {-6,...,6} {
  \pgfmathtruncatemacro{\r}{mod(\i+99,3)}
  \ifnum\r=0\def\col{blue!70!black}\else
    \ifnum\r=1\def\col{red!75!black}\else\def\col{teal!70!black}\fi
  \fi
  \fill[\col] (\i,0) circle (2.7pt);
  \node[font=\scriptsize,below=5pt,\col] at (\i,0) {$\i$};
}
\draw[blue!70!black,thick,rounded corners] (-5.35,0.48)--(-2.65,0.48)
  node[midway,above,font=\small] {$\overline0$};
\draw[red!75!black,thick,rounded corners] (-4.35,1.03)--(-1.65,1.03)
  node[midway,above,font=\small] {$\overline1$};
\draw[teal!70!black,thick,rounded corners] (-3.35,1.58)--(-0.65,1.58)
  node[midway,above,font=\small] {$\overline2$};
\node[font=\small,align=center] at (3.6,1.25)
  {Kolor oznacza klasę\\reszt modulo $3$.};
\end{tikzpicture}
\end{center}

Na rysunku punkty tego samego koloru leżą w jednej klasie równoważności.
Klasa nie jest jedną liczbą: na przykład $\overline1$ zawiera
$\dots,-5,-2,1,4,7,\dots$. W ilorazie traktujemy je jako jeden element.

\subsection{Wielomiany, funkcje i iloczyny}

\begin{przyklad}[Pierścień wielomianów]
Jeżeli $A$ jest pierścieniem, to $A[x]$ oznacza zbiór formalnych sum
$f(x)=a_0+a_1x+\dots+a_dx^d$, gdzie $d\ge0$ i $a_j\in A$, oraz wielomian
zerowy. Dwa wielomiany są równe wtedy i tylko wtedy, gdy mają identyczne
współczynniki przy każdej potędze $x$. Dodajemy je współczynnikami,
a mnożenie określa rozdzielność i reguła $x^ix^j=x^{i+j}$.
Przykładowo w $\Z[x]$ mamy
$(2+x)(1+3x)=2+7x+3x^2$.
Nie utożsamiamy na tym etapie formalnego wielomianu z funkcją:
na $\F_2=\Z/2\Z$ wielomian $x^2-x$ nie jest wielomianem zerowym,
choć w obu punktach ciała $\F_2$ przyjmuje wartość zero.
\end{przyklad}

\begin{uwaga}[Wielomian formalny a funkcja wielomianowa]
\label{uw:formalna-funkcja}
Dla dowolnego pierścienia przemiennego $A$ podstawienie określa
odwzorowanie $A[x]\to A^A$, które wielomianowi $f$ przyporządkowuje
funkcję $a\mapsto f(a)$. Nie zawsze jest ono różnowartościowe.
Przykładowo w $\mathbb F_2[x]$ wielomiany $x^2-x$ oraz $0$ są różne,
ale wyznaczają tę samą funkcję na $\mathbb F_2$, gdyż
$0^2-0=0$ i $1^2-1=0$.

Jeśli natomiast $k$ jest ciałem nieskończonym, to odwzorowanie
$k[x]\to k^k$ jest różnowartościowe: gdy dwa wielomiany wyznaczają tę samą
funkcję, ich różnica znika w każdym punkcie $k$, a niezerowy wielomian
ma tylko skończenie wiele pierwiastków (twierdzenie~\ref{tw:korzenie}).
W przypadku $n$ zmiennych wniosek jest analogiczny. Dowód prowadzimy
indukcyjnie: dla wielomianu
$f(x_1,\dots,x_n)=\sum_{j=0}^d a_j(x_1,\dots,x_{n-1})x_n^j$
zerującego się na całym $k^n$ ustalamy pierwsze $n-1$ współrzędnych.
Otrzymany wielomian jednej zmiennej jest zerowy, więc wszystkie
$a_j$ znikają na $k^{n-1}$; indukcja pokazuje, że wszystkie są
wielomianami zerowymi.

W geometrii algebraicznej rozróżniamy więc 
\emph{pierścień formalnych wielomianów}, 
\emph{funkcje na punktach nad $k$} oraz 
\emph{pierścień zadany relacjami} (iloraz przez ideał).
Te obiekty nie są zamienne bez dodatkowych założeń.
\end{uwaga}

\begin{przyklad}[Pierścień funkcji]
Niech $X$ będzie niepustym zbiorem, a $A$ pierścieniem.
Na zbiorze $A^X$ wszystkich odwzorowań $f\colon X\to A$ definiujemy
$(f+g)(x)=f(x)+g(x)$ oraz $(fg)(x)=f(x)g(x)$.
Zero i jedynka są funkcjami stałymi równymi odpowiednio $0$ i $1$.
Sprawdzenie rozdzielności sprowadza się do równości
$(f(g+h))(x)=f(x)(g(x)+h(x))=f(x)g(x)+f(x)h(x)$
dla każdego $x\in X$. Otrzymujemy pierścień.
\end{przyklad}

\begin{przyklad}[Iloczyn pierścieni]
Dla pierścieni $A,B$ zbiór $A\times B$ jest pierścieniem z działaniami
$(a,b)+(a',b')=(a+a',b+b')$ i
$(a,b)(a',b')=(aa',bb')$. Jedynką jest $(1_A,1_B)$.
W pierścieniu $\R\times\R$ zachodzi $(1,0)(0,1)=(0,0)$,
choć żaden z dwóch czynników nie jest zerowy.
\end{przyklad}

\begin{przyklad}[Pierścień macierzy — przykład ostrzegawczy]
W zbiorze $M_2(\R)$ macierzy $2\times2$ mamy naturalne dodawanie,
mnożenie i macierz jednostkową. Jest to pierścień, ale nieprzemienny:
przyjmując
$P=\begin{psmallmatrix}1&0\\0&0\end{psmallmatrix}$ oraz
$Q=\begin{psmallmatrix}0&1\\0&0\end{psmallmatrix}$,
otrzymujemy $PQ=Q$, podczas gdy $QP=0$.
Od tej chwili będziemy głównie pracować w pierścieniach przemiennych,
ponieważ to właśnie ich spektra stanowią podstawę klasycznej geometrii
algebraicznej.
\end{przyklad}

\section{Elementy odwracalne, dzielniki zera i elementy nilpotentne}

\begin{definicja}[Jednostka, dzielnik zera, element nilpotentny]
Niech $A$ będzie pierścieniem przemiennym.
\begin{enumerate}[label=\arabic*)]
\item Element $a\in A$ jest \emph{odwracalny} (jest \emph{jednostką}),
jeżeli istnieje $b\in A$ takie, że $ab=1$. Wtedy $b$ oznaczamy $a^{-1}$.
Zbiór jednostek oznaczamy $A^\times$.
\item Niezerowy $a\in A$ jest \emph{dzielnikiem zera}, jeżeli istnieje
niezerowy $b\in A$ taki, że $ab=0$.
\item Element $a\in A$ jest \emph{nilpotentny}, jeżeli $a^m=0$
dla pewnego całkowitego $m\ge1$.
\item Element $e\in A$ jest \emph{idempotentny}, jeżeli $e^2=e$.
\end{enumerate}
\end{definicja}

\begin{stwierdzenie}
Niech $A$ będzie pierścieniem przemiennym. Jednostki tworzą grupę
abelową $(A^\times,\cdot)$. Jeżeli $a$ jest jednostką i $ab=ac$,
to $b=c$. Żadna jednostka w niezerowym pierścieniu nie jest dzielnikiem zera.
\end{stwierdzenie}
\begin{proof}
Jedynka jest jednostką, a odwrotność jednostki również jest jednostką.
Jeżeli $ab=1$ i $cd=1$, to $(ac)(bd)=(ab)(cd)=1$, więc iloczyn jednostek
jest jednostką. Łączność i przemienność dziedziczymy z $A$.
Z $ab=ac$ po pomnożeniu przez $a^{-1}$ dostajemy
$(a^{-1}a)b=(a^{-1}a)c$, czyli $b=c$.
Jeśli $ab=0$ i $a$ jest jednostką, to $b=a^{-1}0=0$.
\end{proof}

\begin{przyklad}[Które klasy modulo $n$ są odwracalne?]
W $\Z/12\Z$ jednostkami są dokładnie
$\overline1,\overline5,\overline7,\overline{11}$.
Na przykład $\overline5^2=\overline{25}=\overline1$.
Element $\overline6$ nie jest jednostką, ponieważ
$\overline6\,\overline2=\overline0$ i $\overline2\ne\overline0$.
Poniższe twierdzenie daje ogólną odpowiedź.
\end{przyklad}

\begin{lemat}[Tożsamość Bézouta dla liczb całkowitych]
\label{lem:bezout-int}
Niech $a,b\in\Z$ nie będą jednocześnie zerowe. Wówczas istnieją
$u,v\in\Z$ takie, że $ua+vb=\gcd(a,b)$.
\end{lemat}
\begin{proof}
Rozważmy zbiór $S=\{ua+vb:u,v\in\Z\}$.
Jest on zamknięty na dodawanie i odejmowanie. Ponieważ przynajmniej
jedna z liczb $a,b$ jest niezerowa, $S$ zawiera pewną liczbę dodatnią
(w razie potrzeby zmieniamy znak). Z zasady dobrego uporządkowania
liczb naturalnych istnieje najmniejsza dodatnia liczba $d\in S$.
Z definicji $S$ mamy $d=ua+vb$ dla pewnych $u,v\in\Z$.

Dzieląc $a$ przez $d$ z resztą, otrzymujemy $a=qd+r$,
gdzie $0\le r<d$. Ponieważ $a,d\in S$, także $r=a-qd\in S$.
Minimalność $d$ wymusza $r=0$, a zatem $d\mid a$.
Identycznie dowodzimy $d\mid b$.

Jeśli $c$ jest dowolnym wspólnym dzielnikiem $a$ i $b$,
to $c$ dzieli również ich kombinację $ua+vb=d$.
Wykazaliśmy, że $d$ jest dodatnim wspólnym dzielnikiem,
przez który dzieli się każdy wspólny dzielnik $a,b$.
Jest to właśnie największy wspólny dzielnik: $d=\gcd(a,b)$.
\end{proof}

\begin{twierdzenie}[Jednostki modulo $n$]
\label{tw:jednostki-mod}
Niech $n\ge2$ oraz $a\in\Z$. Klasa $\overline a\in\Z/n\Z$
jest jednostką wtedy i tylko wtedy, gdy $\gcd(a,n)=1$.
\end{twierdzenie}
\begin{proof}
Jeśli $\overline a$ jest jednostką, to istnieje $b\in\Z$ takie, że
$ab\equiv1\pmod n$. Oznacza to, że $ab-1=kn$ dla pewnego $k\in\Z$,
czyli $ab-kn=1$. Każdy wspólny dzielnik $a$ i $n$ dzieli lewą stronę,
a więc dzieli $1$. Stąd $\gcd(a,n)=1$.

Na odwrót, dla względnie pierwszych $a,n$ lemat~\ref{lem:bezout-int} daje $u,v\in\Z$ takie, że $au+nv=1$.
Przechodząc do klas modulo $n$, otrzymujemy
$\overline a\,\overline u=\overline1$. Zatem $\overline a$ jest jednostką.
\end{proof}

\begin{uwaga}[Skąd bierze się tożsamość Bézouta?]
Można ją otrzymać algorytmem Euklidesa. Dla $a=5,n=12$ mamy
$12=2\cdot5+2$, $5=2\cdot2+1$, a z podstawienia
$1=5-2(12-2\cdot5)=5\cdot5-2\cdot12$.
Stąd odwrotnością $\overline5$ w $\Z/12\Z$ jest $\overline5$.
Analogiczny algorytm dla wielomianów pojawi się dalej.
\end{uwaga}

\begin{definicja}[Dziedzina całkowitości i ciało]
\label{def:dziedzina-cialo}
\emph{Dziedziną całkowitości} nazywamy niezerowy pierścień przemienny,
w którym z $ab=0$ wynika $a=0$ lub $b=0$.
\emph{Ciałem} nazywamy niezerowy pierścień przemienny,
w którym każdy element różny od zera jest jednostką.
\end{definicja}

Każde ciało jest dziedziną: jeśli $ab=0$ i $a\ne0$, możemy pomnożyć
przez $a^{-1}$ i otrzymać $b=0$. Przykładami ciał są $\Q,\R,\C$,
a przykładem dziedziny niebędącej ciałem jest $\Z$.

\begin{stwierdzenie}[Skracanie w dziedzinie]
Dla niezerowego pierścienia przemiennego $A$ następujące warunki są
równoważne: (i) $A$ jest dziedziną całkowitości;
(ii) dla wszystkich $a,b,c\in A$ z $a\ne0$ i $ab=ac$ wynika $b=c$.
\end{stwierdzenie}
\begin{proof}
Jeśli $A$ jest dziedziną, to $ab=ac$ daje $a(b-c)=0$,
a $a\ne0$ wymusza $b-c=0$. Jeśli zachodzi (ii) oraz $ab=0$ przy $a\ne0$,
to $ab=a0$, więc $b=0$. Jest to definicja dziedziny.
\end{proof}

\begin{twierdzenie}[Skończona dziedzina jest ciałem]
\label{tw:skonczona}
Każda skończona dziedzina całkowitości $A$ jest ciałem.
\end{twierdzenie}
\begin{proof}
Ustalmy $a\in A$ z $a\ne0$ i rozważmy odwzorowanie
$m_a\colon A\to A$ zadane przez $m_a(x)=ax$.
Jeżeli $m_a(x)=m_a(y)$, to $ax=ay$. W dziedzinie możemy skrócić
niezerowy czynnik $a$, więc $x=y$; odwzorowanie $m_a$ jest różnowartościowe.
Ponieważ $A$ jest skończony, każde różnowartościowe odwzorowanie $A\to A$
jest też na. Zatem $1=m_a(b)=ab$ dla pewnego $b\in A$.
Dowolny element niezerowy $a$ ma odwrotność, a więc $A$ jest ciałem.
\end{proof}

\begin{wniosek}
Jeżeli $p$ jest liczbą pierwszą, to $\Z/p\Z$ jest ciałem.
Jeżeli $n\ge2$ jest złożone, to $\Z/n\Z$ nie jest dziedziną.
\end{wniosek}
\begin{proof}
Dla pierwszego $p$ każda liczba $1\le a\le p-1$ spełnia
$\gcd(a,p)=1$, więc każda niezerowa klasa jest odwracalna
na mocy twierdzenia~\ref{tw:jednostki-mod}.
Jeśli $n=rs$ z $1<r,s<n$, to
$\overline r\ne0$ i $\overline s\ne0$, ale
$\overline r\,\overline s=\overline n=0$.
\end{proof}

\begin{przyklad}[Nilpotentność nie jest tym samym co bycie dzielnikiem zera]
W $\Z/12\Z$ element $\overline6$ jest nilpotentny,
gdyż $\overline6^2=\overline{36}=0$.
Natomiast $\overline4$ jest idempotentny:
$\overline4^2=\overline{16}=\overline4$.
Element $\overline4$ jest też dzielnikiem zera, bo
$\overline4\,\overline3=0$, ale nie jest nilpotentny:
z $e^2=e$ wynika indukcyjnie $e^m=e$ dla każdego $m\ge1$.
W niezerowej dziedzinie nie ma niezerowych elementów nilpotentnych
ani idempotentów innych niż $0$ i $1$.
\end{przyklad}

\begin{lemat}[Nilpotentność i odwracalność]
Niech $A$ będzie pierścieniem przemiennym i niech $u\in A$ spełnia
$u^m=0$ dla pewnego $m\ge1$. Wtedy $1-u$ jest jednostką, a jej
odwrotnością jest $1+u+\dots+u^{m-1}$.
\end{lemat}
\begin{proof}
Korzystamy ze wzoru na sumę geometryczną, który działa w każdym
pierścieniu przemiennym:
$(1-u)(1+u+\dots+u^{m-1})=1-u^m=1$.
\end{proof}


\subsection{Charakterystyka pierścienia}

\begin{definicja}[Charakterystyka]
Niech $A$ będzie niezerowym pierścieniem.
Jeżeli istnieje dodatnia liczba całkowita $n$ spełniająca
$\underbrace{1+\dots+1}_{n\ \text{składników}}=0$,
to najmniejszą taką liczbę nazywamy \emph{charakterystyką} $A$
i oznaczamy $\operatorname{char}A$.
Jeżeli nie ma takiej liczby, przyjmujemy $\operatorname{char}A=0$.
\end{definicja}

\begin{przyklad}
Mamy $\operatorname{char}\Z=\operatorname{char}\Q
=\operatorname{char}\R=\operatorname{char}\C=0$,
a dla $n\ge2$ zachodzi $\operatorname{char}(\Z/n\Z)=n$.
Dla dowolnego niezerowego pierścienia $A$ charakterystyka $A[x]$ jest równa
charakterystyce $A$, ponieważ $1$ w $A[x]$ jest wielomianem stałym $1$.
\end{przyklad}

\begin{stwierdzenie}
Niech $A$ będzie dziedziną całkowitości. Wówczas
$\operatorname{char}A=0$ albo $\operatorname{char}A=p$
dla pewnej liczby pierwszej $p$.
\end{stwierdzenie}
\begin{proof}
Przypuśćmy, że charakterystyka jest dodatnia i wynosi $n$.
Gdyby $n=rs$ dla $1<r,s<n$, to
$(\underbrace{1+\dots+1}_{r})(\underbrace{1+\dots+1}_{s})
=\underbrace{1+\dots+1}_{rs}=0$.
Z minimalności $n$ oba czynniki są niezerowe, co przeczy
brakowi dzielników zera w dziedzinie. Zatem $n$ jest liczbą pierwszą.
\end{proof}

\section{Homomorfizmy: kiedy dwa opisy zachowują rachunek?}

\begin{definicja}[Homomorfizm pierścieni]
Niech $A,B$ będą pierścieniami. Odwzorowanie
$\varphi\colon A\to B$ jest \emph{homomorfizmem pierścieni}, jeżeli dla
wszystkich $a,b\in A$ zachodzą równości
$\varphi(a+b)=\varphi(a)+\varphi(b)$,
$\varphi(ab)=\varphi(a)\varphi(b)$ oraz $\varphi(1_A)=1_B$.
Bijektywny homomorfizm nazywamy \emph{izomorfizmem}.
\end{definicja}

Warunku zachowywania jedynki nie pomijamy. Jest on istotny:
zerowe odwzorowanie między niezerowymi pierścieniami zachowuje dodawanie
i mnożenie, ale nie wysyła jedynki na jedynkę.

\begin{stwierdzenie}
Jeżeli $\varphi\colon A\to B$ jest homomorfizmem pierścieni, to
$\varphi(0_A)=0_B$ i $\varphi(-a)=-\varphi(a)$ dla każdego $a\in A$.
Ponadto obraz jednostki $a\in A^\times$ jest jednostką w $B$,
a $\varphi(a^{-1})=\varphi(a)^{-1}$.
\end{stwierdzenie}
\begin{proof}
Ponieważ $\varphi(0_A)=\varphi(0_A+0_A)
=\varphi(0_A)+\varphi(0_A)$, możemy odjąć $\varphi(0_A)$
i otrzymać $\varphi(0_A)=0_B$.
Następnie $\varphi(a)+\varphi(-a)=\varphi(a-a)=0_B$,
więc $\varphi(-a)=-\varphi(a)$.
Gdy $ab=1_A$, to $\varphi(a)\varphi(b)=\varphi(1_A)=1_B$,
co dowodzi ostatniej części.
\end{proof}

\begin{definicja}[Jądro i obraz]
Dla homomorfizmu $\varphi\colon A\to B$ definiujemy
$\kerf\varphi=\{a\in A:\varphi(a)=0_B\}$ oraz
$\im\varphi=\{\varphi(a):a\in A\}\subseteq B$.
\end{definicja}

\begin{przyklad}[Redukcja modulo $n$]
Dla $n\ge2$ odwzorowanie
$\pi_n\colon\Z\to\Z/n\Z$, $a\mapsto\overline a$, jest
surjektywnym homomorfizmem. Jego jądro składa się dokładnie z liczb
podzielnych przez $n$: $\kerf\pi_n=n\Z$.
Odwzorowanie to \emph{zapomina}, którą konkretną liczbą całkowitą
reprezentowana jest klasa, ale zachowuje dodawanie i mnożenie.
\end{przyklad}

\begin{przyklad}[Podstawienie do wielomianu]
Niech $A$ będzie pierścieniem i $a\in A$. Odwzorowanie
$\ev_a\colon A[x]\to A$ dane przez $\ev_a(f)=f(a)$
jest homomorfizmem pierścieni. Dla $f=\sum_i f_ix^i$ i
$g=\sum_jg_jx^j$ mamy
$\ev_a(fg)=\sum_{i,j}f_ig_ja^{i+j}
=(\sum_if_ia^i)(\sum_jg_ja^j)=f(a)g(a)$.
Dodawanie i zachowywanie jedynki sprawdzamy bezpośrednio.
Surjektywność wynika z tego, że każda stała $b\in A$ jest wartością
stałego wielomianu $b$.
\end{przyklad}

\begin{uwaga}[Dlaczego jądro jest ważniejsze od samego wzoru?]
Dwa elementy $a,b\in A$ mają ten sam obraz wtedy i tylko wtedy, gdy
$\varphi(a-b)=0$, czyli $a-b\in\kerf\varphi$.
Jądro dokładnie opisuje więc \emph{jakie elementy stają się
nierozróżnialne} po zastosowaniu homomorfizmu. Aby z takiego
utożsamiania otrzymać nowy pierścień, potrzebujemy pojęcia ideału.
\end{uwaga}

\section{Ideały: podzbiory nadające się do ilorazowania}

\begin{definicja}[Ideał]
Niech $A$ będzie pierścieniem przemiennym. Podzbiór $I\subseteq A$
nazywamy \emph{ideałem}, jeżeli
\begin{enumerate}[label=(I\arabic*)]
\item $0\in I$ i $a-b\in I$ dla wszystkich $a,b\in I$;
\item $ra\in I$ dla wszystkich $r\in A$ i $a\in I$.
\end{enumerate}
Ideał jest \emph{właściwy}, gdy $I\ne A$.
\end{definicja}

Pierwszy warunek znaczy, że $(I,+)$ jest podgrupą $(A,+)$.
Drugi nazywa się \emph{pochłanianiem przez mnożenie}.
Można wyobrażać sobie ideał jako zbiór elementów, które chcemy
\emph{uznać za zero}: jeśli $a$ ma stać się zerem, to również
$ra$ musi stać się zerem, aby mnożenie pozostawało sensowne.

\begin{przyklad}[Dlaczego $2\Z$ jest ideałem]
Suma i różnica dwóch liczb parzystych są parzyste,
a iloczyn dowolnej liczby całkowitej z parzystą jest parzysty.
Zatem $2\Z$ jest ideałem $\Z$. Nie jest jednak podpierścieniem
w naszym znaczeniu, ponieważ $1\notin2\Z$.
\end{przyklad}

\begin{stwierdzenie}[Test właściwości ideału]
Niech $I$ będzie ideałem pierścienia $A$. Wówczas
$I=A$ wtedy i tylko wtedy, gdy $1\in I$.
Jeśli $I$ zawiera jednostkę, to $I=A$.
\end{stwierdzenie}
\begin{proof}
Jeżeli $I=A$, to oczywiście $1\in I$. Jeśli zaś $1\in I$,
to dla każdego $a\in A$ mamy $a=a\cdot1\in I$, więc $I=A$.
Gdy $u\in I$ jest jednostką, pochłanianie daje
$1=u^{-1}u\in I$, więc z pierwszej części $I=A$.
\end{proof}

\begin{definicja}[Ideał generowany przez zbiór]
Dla podzbioru $S\subseteq A$ przez $(S)$ oznaczamy najmniejszy ideał
pierścienia $A$ zawierający $S$, to znaczy przecięcie wszystkich ideałów
zawierających $S$. Jeśli $S=\{a_1,\dots,a_r\}$, piszemy
$(a_1,\dots,a_r)$. Ideał $(a)$ nazywamy \emph{głównym}.
\end{definicja}

\begin{stwierdzenie}[Jawny opis generowania]
Dla $S\subseteq A$ ideał $(S)$ składa się z wszystkich skończonych sum
$\sum_{j=1}^m r_js_j$, gdzie $m\ge0$, $r_j\in A$ i $s_j\in S$.
W szczególności $(a_1,\dots,a_r)=\{r_1a_1+\dots+r_ra_r:r_j\in A\}$
oraz $(a)=\{ra:r\in A\}$.
\end{stwierdzenie}
\begin{proof}
Nazwijmy zbiór wszystkich wskazanych sum $J$.
Suma pusta daje $0\in J$. Suma oraz różnica dwóch takich wyrażeń
też jest sumą odpowiedniej postaci (minus przenosimy na współczynniki).
Dla $t\in A$ mamy $t\sum_jr_js_j=\sum_j(tr_j)s_j\in J$.
Zatem $J$ jest ideałem. Ponadto każdy $s\in S$ należy do $J$,
bo $s=1\cdot s$.
Jeśli $I$ jest dowolnym ideałem zawierającym $S$, to pochłanianie
daje $r_js_j\in I$ dla wszystkich $j$, a zamknięcie na dodawanie
daje $\sum_jr_js_j\in I$. Stąd $J\subseteq I$.
Zbiór $J$ jest więc najmniejszym ideałem zawierającym $S$, czyli $J=(S)$.
\end{proof}

\begin{przyklad}[Dwa generatory mogą być potrzebne]
W $\Z$ mamy $(6,10)=(2)$: każda kombinacja $6a+10b$ jest parzysta,
a z drugiej strony $2=2\cdot6-1\cdot10$.
Natomiast w $k[x,y]$, gdzie $k$ jest ciałem, ideał $(x,y)$
nie jest generowany przez jeden wielomian.
Dowód tego drugiego faktu podamy po omówieniu dzielenia wielomianów;
ważne jest na razie rozróżnienie ideałów jedno- i wielogeneratorowych.
\end{przyklad}

\begin{stwierdzenie}[Suma, przecięcie i iloczyn ideałów]
Dla ideałów $I,J$ pierścienia $A$ zbiory
$I\cap J$ oraz $I+J=\{a+b:a\in I,\ b\in J\}$ są ideałami.
Suma $I+J$ jest najmniejszym ideałem zawierającym $I\cup J$.
Zbiór
$IJ=\{\sum_{r=1}^m a_rb_r:m\ge0,\ a_r\in I,\ b_r\in J\}$
także jest ideałem, a ponadto $IJ\subseteq I\cap J$.
\end{stwierdzenie}
\begin{proof}
Przecięcie dwóch ideałów zawiera zero i jest zamknięte na odejmowanie
oraz mnożenie przez elementy $A$, bo każda z tych własności zachodzi
w obu ideałach. Dla sumy: $(a+b)-(a'+b')=(a-a')+(b-b')\in I+J$,
a $r(a+b)=ra+rb\in I+J$. Ideał $I+J$ zawiera $I$ i $J$,
bo $0$ należy do obu; dowolny ideał zawierający $I$ i $J$ zawiera
również wszystkie sumy $a+b$.

Zbiór $IJ$ jest ideałem: można łączyć i odejmować skończone sumy,
a mnożenie przez $r\in A$ daje
$r\sum a_rb_r=\sum(ra_r)b_r$ z $ra_r\in I$.
Każdy składnik $a_rb_r$ należy do $I$ (bo $a_r\in I$)
i do $J$ (bo $b_r\in J$), a więc $IJ\subseteq I\cap J$.
\end{proof}

\begin{przyklad}[Przecięcie nie musi równać się iloczynowi]
W $\Z$ mamy $(2)\cap(4)=(4)$, natomiast $(2)(4)=(8)$.
Inaczej jest dla $(2)$ i $(3)$: wtedy
$(2)\cap(3)=(6)=(2)(3)$. Równość w drugim przypadku wiąże się
z faktem, że $2$ i $3$ są względnie pierwsze.
\end{przyklad}

\begin{definicja}[Ideały współmaksymalne]
Ideały $I,J\subseteq A$ nazywamy \emph{współmaksymalnymi}, jeśli
$I+J=A$. Równoważnie istnieją $a\in I$, $b\in J$ takie, że $a+b=1$.
\end{definicja}

\begin{stwierdzenie}
Niech $A$ będzie pierścieniem przemiennym z jedynką, a $I,J\subseteq A$ ideałami.
Jeżeli $I+J=A$, to $I\cap J=IJ$.
\end{stwierdzenie}
\begin{proof}
Wiemy już, że $IJ\subseteq I\cap J$.
Wybierzmy $a\in I,b\in J$ z $a+b=1$. Jeśli $c\in I\cap J$,
to $c=c(a+b)=ca+cb$. Ponieważ $c\in J$ i $a\in I$,
mamy $ca\in IJ$; ponieważ $c\in I$ i $b\in J$,
także $cb\in IJ$. Zatem $c\in IJ$, co daje drugą inkluzję.
\end{proof}

\section{Pierścienie ilorazowe: co oznacza utożsamienie elementów?}

\begin{definicja}[Pierścień ilorazowy]
Niech $I$ będzie ideałem pierścienia $A$. Definiujemy relację
$a\sim_Ib$ wtedy i tylko wtedy, gdy $a-b\in I$.
Klasę elementu $a$ oznaczamy $a+I=\{a+i:i\in I\}$.
Zbiór tych klas zapisujemy $A/I$.
Działania na klasach określamy wzorami
$(a+I)+(b+I)=(a+b)+I$ oraz $(a+I)(b+I)=ab+I$.
\end{definicja}

\begin{twierdzenie}[Poprawność konstrukcji ilorazowej]
\label{tw:iloraz}
Dla każdego ideału $I\subseteq A$ relacja $\sim_I$ jest relacją
równoważności, powyższe działania są dobrze określone,
a $A/I$ jest pierścieniem przemiennym z zerem $0+I$
i jedynką $1+I$.
\end{twierdzenie}
\begin{proof}
Zwrotność: $a-a=0\in I$.
Symetria: jeśli $a-b\in I$, to $b-a=-(a-b)\in I$.
Przechodniość: jeśli $a-b,b-c\in I$, to
$a-c=(a-b)+(b-c)\in I$.

Załóżmy teraz, że $a+I=a'+I$ i $b+I=b'+I$.
Oznacza to $a-a'\in I$ i $b-b'\in I$.
Dla dodawania mamy
$(a+b)-(a'+b')=(a-a')+(b-b')\in I$.
Dla mnożenia korzystamy z rozkładu
$ab-a'b'=(a-a')b+a'(b-b')$.
Oba składniki po prawej należą do $I$ na mocy własności pochłaniania,
więc $ab+I=a'b'+I$. To jest kluczowy punkt dowodu:
\emph{bez pochłaniania mnożenie klas nie musiałoby być dobrze określone}.

Aksjomaty pierścienia wynikają z aksjomatów w $A$;
na przykład łączność mnożenia to
$((a+I)(b+I))(c+I)=(ab)c+I=a(bc)+I
=(a+I)((b+I)(c+I))$.
Analogicznie sprawdzamy rozdzielność, przemienność i istnienie
odwrotności addytywnej $(-a)+I$.
\end{proof}

\begin{przyklad}[Dlaczego podgrupa addytywna nie wystarczy]
W $A=\Z[x]$ zbiór $H=\Z$ jest podgrupą grupy addytywnej $(A,+)$,
ale nie jest ideałem $A$, ponieważ $1\in H$ i $x\cdot1=x\notin H$.
Gdybyśmy w klasach addytywnych modulo $H$ chcieli mnożyć,
to $0$ i $1$ byłyby w tej samej klasie, ale ich iloczyny przez $x$
różniłyby się o $x$, który nie należy do $H$.
Wynik mnożenia zależałby zatem od reprezentanta.
\end{przyklad}

\begin{definicja}[Rzutowanie kanoniczne]
Dla ideału $I\subseteq A$ odwzorowanie
$\pi_I\colon A\to A/I$, $\pi_I(a)=a+I$,
nazywamy \emph{rzutowaniem kanonicznym}.
\end{definicja}

\begin{stwierdzenie}
Niech $A$ będzie pierścieniem przemiennym z jedynką, a $I\subseteq A$
ideałem. Rzutowanie kanoniczne $\pi_I:A\to A/I$ jest
surjektywnym homomorfizmem pierścieni i $\kerf\pi_I=I$.
\end{stwierdzenie}
\begin{proof}
Każda klasa w $A/I$ jest obrazem swojego reprezentanta,
więc $\pi_I$ jest na. Zachowywanie działań wynika z ich definicji.
Ponadto $\pi_I(a)=0+I$ jest równoważne $a-0\in I$,
a więc $a\in I$.
\end{proof}

\begin{twierdzenie}[Pierwsze twierdzenie o izomorfizmie]
\label{tw:izo}
Dla każdego homomorfizmu pierścieni $\varphi\colon A\to B$
istnieje dokładnie jeden izomorfizm
$\widetilde\varphi\colon A/\kerf\varphi\to\im\varphi$
spełniający $\widetilde\varphi(a+\kerf\varphi)=\varphi(a)$
dla wszystkich $a\in A$.
\end{twierdzenie}
\begin{proof}
Najpierw sprawdzamy, że $\kerf\varphi$ jest ideałem:
$\varphi(0)=0$; jeśli $a,b$ należą do jądra, to
$\varphi(a-b)=\varphi(a)-\varphi(b)=0$;
jeśli $r\in A$ i $a\in\kerf\varphi$, to
$\varphi(ra)=\varphi(r)\varphi(a)=0$.

Definiujemy $\widetilde\varphi(a+\kerf\varphi)=\varphi(a)$.
Jeśli $a+\kerf\varphi=b+\kerf\varphi$, to
$a-b\in\kerf\varphi$, zatem $\varphi(a)=\varphi(b)$.
Wzór nie zależy więc od reprezentanta. Homomorfizm zachowuje działania,
bo na przykład
$\widetilde\varphi((a+K)(b+K))=\varphi(ab)
=\varphi(a)\varphi(b)$ dla $K=\kerf\varphi$.
Z definicji obrazu jest surjektywny na $\im\varphi$.
Jeżeli $\widetilde\varphi(a+K)=\widetilde\varphi(b+K)$,
to $\varphi(a-b)=0$, a więc $a-b\in K$, czyli $a+K=b+K$.
Jest zatem różnowartościowy. Jedyność wynika z tego, że każda klasa
w ilorazie jest postaci $a+K$, na której żądany wzór ustala wartość.
\end{proof}

\begin{center}
\begin{tikzpicture}[>=Stealth,node distance=2.15cm]
\node (a) {$A$};
\node[right=3.0cm of a] (im) {$\im\varphi$};
\node[below=1.6cm of a] (q) {$A/\kerf\varphi$};
\draw[->] (a)--node[above] {$\varphi$} (im);
\draw[->] (a)--node[left] {$\pi$} (q);
\draw[->] (q)--node[below right] {$\widetilde\varphi\ \cong$} (im);
\end{tikzpicture}
\end{center}
W diagramie obie drogi z $A$ do $\im\varphi$ dają ten sam wynik.
Mówimy, że diagram jest \emph{przemienny}.

\begin{przyklad}[Wartość wielomianu i iloraz]
Dla $a\in A$ homomorfizm $\ev_a\colon A[x]\to A$ jest surjektywny.
Udowodnimy teraz, że $\kerf\ev_a=(x-a)$.
Dla dowolnego $f(x)=\sum_{j=0}^m c_jx^j$ zapisujemy
\[
 f(x)-f(a)=\sum_{j=1}^m c_j(x^j-a^j)
 =(x-a)\sum_{j=1}^m c_j\sum_{r=0}^{j-1}x^{j-1-r}a^r.
\]
Jeżeli $f(a)=0$, to $f(x)$ jest wielokrotnością $x-a$.
Jeśli odwrotnie $f(x)=(x-a)g(x)$, to podstawienie $x=a$ daje $f(a)=0$.
Wykorzystując twierdzenie~\ref{tw:izo}, otrzymujemy
$A[x]/(x-a)\cong A$.
To elementarna wersja zasady: \emph{równanie $x=a$ pozwala usunąć
jedną zmienną}.
\end{przyklad}

\subsection{Ideały w ilorazie i ich odpowiedniki w pierścieniu}

\begin{twierdzenie}[Odpowiedniość ideałów]
\label{tw:odp-idealow}
Niech $A$ będzie pierścieniem i $I$ jego ideałem.
Przyporządkowanie $J\mapsto J/I=\{a+I:a\in J\}$ jest bijekcją
między ideałami $J\subseteq A$ spełniającymi $I\subseteq J$
a ideałami pierścienia $A/I$. Odwrotne przyporządkowanie ma postać
$K\mapsto\pi_I^{-1}(K)$. Bijekcja zachowuje inkluzje.
\end{twierdzenie}
\begin{proof}
Dla $I\subseteq J$ zbiór $J/I$ jest ideałem $A/I$:
zerowa klasa do niego należy, różnice klas elementów $J$ pozostają
w $J/I$, a dla $r+I\in A/I$ i $a+I\in J/I$ iloczyn
$ra+I$ leży w $J/I$ dzięki $ra\in J$.

Jeśli $K$ jest ideałem $A/I$, to przeciwobraz $\pi_I^{-1}(K)$ jest
ideałem $A$ (sprawdzamy warunki ideału, przenosząc je przez homomorfizm).
Zawiera $I$, ponieważ $I=\kerf\pi_I$, a zero należy do $K$.

Dla $I\subseteq J$ zachodzi
$\pi_I^{-1}(J/I)=J$: jeśli $a+I=j+I$ dla pewnego $j\in J$,
to $a-j\in I\subseteq J$, a więc $a\in J$; odwrotna inkluzja
jest oczywista. Z kolei dla ideału $K\subseteq A/I$ mamy
$\pi_I(\pi_I^{-1}(K))=K$ dzięki surjektywności $\pi_I$.
To dowodzi, że przyporządkowania są odwrotne.
Jeżeli $J_1\subseteq J_2$, to oczywiście $J_1/I\subseteq J_2/I$;
analogicznie dla przeciwobrazów.
\end{proof}

\begin{przyklad}[Wszystkie ideały $\Z/12\Z$]
Za chwilę pokażemy, że wszystkie ideały $\Z$ mają postać $d\Z$.
Te spośród nich, które zawierają $12\Z$, są generowane przez
dzielniki $12$ (z dokładnością do znaku):
$12\Z,6\Z,4\Z,3\Z,2\Z,\Z$.
Z twierdzenia~\ref{tw:odp-idealow} otrzymujemy sześć ideałów
pierścienia $\Z/12\Z$, oznaczanych odpowiednio
$(0),(\overline6),(\overline4),(\overline3),(\overline2),(1)$.
Ich uporządkowanie przez zawieranie można narysować następująco:
\end{przyklad}

\begin{center}
\begin{tikzpicture}[>=Stealth, every node/.style={font=\small,inner sep=3pt}]
\node (top) at (0,3.5) {$(1)$};
\node (two) at (-1.2,2.55) {$(\overline2)$};
\node (three) at (1.2,2.55) {$(\overline3)$};
\node (four) at (-1.2,1.55) {$(\overline4)$};
\node (six) at (1.2,1.55) {$(\overline6)$};
\node (zero) at (0,0.55) {$(0)$};
\draw (top)--(two)--(four)--(zero);
\draw (top)--(three)--(six)--(zero);
\draw (two)--(six);
\end{tikzpicture}
\end{center}
Każdy odcinek oznacza relację bezpośredniego zawierania:
ideał na dole jest zawarty w ideale na górze. Na przykład
$(\overline6)\subseteq(\overline2)$ i
$(\overline6)\subseteq(\overline3)$, ponieważ $6$ jest wielokrotnością
zarówno $2$, jak i $3$. Rysunek nie jest wykresem funkcji,
lecz diagramem porządku częściowego.


\subsection{Chińskie twierdzenie o resztach}

Istnieje ważna sytuacja, w której pierścień ilorazowy można rozłożyć
na iloczyn dwóch prostszych pierścieni. Pozwala to rozwiązywać układy
kongruencji, a zarazem tłumaczy obecność idempotentów w pierścieniach.

\begin{twierdzenie}[Chińskie twierdzenie o resztach]
\label{tw:crt}
Niech $A$ będzie pierścieniem przemiennym, a $I,J$ jego ideałami
takimi, że $I+J=A$. Wtedy odwzorowanie
\[
 \Theta\colon A/(I\cap J)\longrightarrow A/I\times A/J,\qquad
 a+(I\cap J)\longmapsto(a+I,a+J),
\]
jest izomorfizmem pierścieni. Ponieważ dla ideałów współmaksymalnych
$I\cap J=IJ$, można równoważnie zapisać $A/IJ\cong A/I\times A/J$.
\end{twierdzenie}
\begin{proof}
Rozważmy najpierw odwzorowanie
$\varphi\colon A\to A/I\times A/J$ zadane wzorem
$\varphi(a)=(a+I,a+J)$. Jest ono homomorfizmem, ponieważ każde
z rzutowań jest homomorfizmem i działania w iloczynie określamy
współrzędnie. Jego jądro składa się z elementów spełniających
jednocześnie $a+I=0+I$ oraz $a+J=0+J$; jest więc równe $I\cap J$.

Pozostaje wykazać surjektywność. Z założenia $I+J=A$ możemy wybrać
$i\in I$ oraz $j\in J$ z $i+j=1$.
Niech $(u+I,v+J)$ będzie dowolnym elementem $A/I\times A/J$.
Weźmy $a=uj+vi$. Ponieważ $j=1-i$, mamy
$a-u=u(j-1)+vi=-ui+vi\in I$.
Analogicznie $i=1-j$, więc
$a-v=uj+v(i-1)=uj-vj\in J$.
Stąd $\varphi(a)=(u+I,v+J)$, co dowodzi, że $\varphi$ jest na.
Pierwsze twierdzenie o izomorfizmie daje
$A/(I\cap J)\cong A/I\times A/J$ i dokładnie wskazany wzór na izomorfizm.
\end{proof}

\begin{przyklad}[Rozkład $\Z/12\Z$]
Ideały $3\Z$ oraz $4\Z$ są współmaksymalne, gdyż
$1=4-3\in3\Z+4\Z$. Ich przecięcie jest równe $12\Z$:
liczba całkowita jest podzielna jednocześnie przez $3$ i przez $4$
dokładnie wtedy, gdy jest podzielna przez $12$.
Chińskie twierdzenie o resztach daje
$\Z/12\Z\cong\Z/3\Z\times\Z/4\Z$.
Izomorfizm wysyła $\overline a$ na
$(a\bmod 3,a\bmod 4)$.
W szczególności $\overline4\mapsto(\overline1,\overline0)$
oraz $\overline9\mapsto(\overline0,\overline1)$.
Tłumaczy to rachunki $\overline4^2=\overline4$ i
$\overline9^2=\overline9$: odpowiadają one oczywistym równościom
$(1,0)^2=(1,0)$ oraz $(0,1)^2=(0,1)$ w iloczynie pierścieni.
Zauważmy także, że $\overline4+\overline9=\overline1$.
\end{przyklad}

\begin{uwaga}[Idea geometryczna]
Idempotenty w pierścieniu pozwalają rozdzielać jego strukturę na
części. W dalszej geometrii algebraicznej odpowiada to rozkładowi
obiektów na pewne rozłączne składowe. Nie należy jednak mylić tego
z rozkładem na składowe nierozkładalne: dwie proste przecinające się
w jednym punkcie tworzą zbiór rozkładalny, ale nierozłączny.
\end{uwaga}

\section{Ideały pierwsze i maksymalne}

\begin{definicja}[Ideał pierwszy]
Właściwy ideał $\mathfrak p\subsetneq A$ jest \emph{pierwszy}, gdy
dla wszystkich $a,b\in A$ zachodzi
$ab\in\mathfrak p\Rightarrow a\in\mathfrak p\ \text{lub}\ b\in\mathfrak p$.
\end{definicja}

\begin{definicja}[Ideał maksymalny]
Właściwy ideał $\mathfrak m\subsetneq A$ jest \emph{maksymalny},
jeżeli dla każdego ideału $J\subseteq A$ z
$\mathfrak m\subseteq J\subseteq A$ zachodzi
$J=\mathfrak m$ lub $J=A$.
\end{definicja}

\begin{uwaga}[Dwa różne znaczenia słowa „pierwszy”]
Element pierwszy i ideał pierwszy to pojęcia pokrewne, ale nie tożsame.
Najpierw definiujemy \emph{ideały pierwsze}, ponieważ działają one
w każdym pierścieniu przemiennym. Elementy pierwsze i nierozkładalne
pojawią się później wraz z teorią podzielności.
\end{uwaga}

\begin{twierdzenie}[Kryterium ilorazowe]
\label{tw:kryterium}
Niech $I$ będzie właściwym ideałem pierścienia $A$.
Wówczas:
\begin{enumerate}[label=\arabic*)]
\item $I$ jest pierwszy wtedy i tylko wtedy, gdy $A/I$ jest dziedziną całkowitości;
\item $I$ jest maksymalny wtedy i tylko wtedy, gdy $A/I$ jest ciałem.
\end{enumerate}
\end{twierdzenie}
\begin{proof}
Ponieważ $I$ jest właściwy, mamy $1+I\ne0+I$.

(1) Jeśli $I$ jest pierwszy i $(a+I)(b+I)=0+I$,
to $ab\in I$, więc $a\in I$ lub $b\in I$.
W takim razie $a+I=0+I$ lub $b+I=0+I$, czyli iloraz jest dziedziną.
Odwrotnie, jeśli $A/I$ jest dziedziną i $ab\in I$, to
$(a+I)(b+I)=0+I$, więc jeden z czynników jest zerową klasą,
co znaczy $a\in I$ lub $b\in I$.

(2) Załóżmy, że $I$ jest maksymalny. Weźmy niezerową klasę $a+I$,
czyli $a\notin I$. Ideał $I+(a)$ zawiera $I$ ściśle,
zatem z maksymalności jest równy $A$.
Istnieją więc $i\in I$ i $r\in A$ takie, że $i+ra=1$.
Przechodząc do ilorazu, dostajemy
$(r+I)(a+I)=1+I$, czyli każda niezerowa klasa ma odwrotność.
Zatem $A/I$ jest ciałem.

Odwrotnie, załóżmy, że $A/I$ jest ciałem.
Niech $J$ będzie ideałem zawierającym $I$ ściśle.
Wybierzmy $a\in J\setminus I$. Klasa $a+I$ jest niezerowa,
więc ma odwrotność $r+I$. Zatem $ra+I=1+I$, czyli $1-ra\in I\subseteq J$.
Ponieważ także $ra\in J$, otrzymujemy $1\in J$, a więc $J=A$.
Właśnie to oznacza maksymalność $I$.
\end{proof}

\begin{wniosek}
Każdy ideał maksymalny pierścienia przemiennego jest pierwszy.
\end{wniosek}
\begin{proof}
Jeśli $\mathfrak m$ jest maksymalny, to $A/\mathfrak m$ jest ciałem,
a więc dziedziną. Z pierwszej części twierdzenia~\ref{tw:kryterium}
wynika, że $\mathfrak m$ jest pierwszy.
\end{proof}

\begin{przyklad}[Ideały w $\Z$]
Pokażemy, że każdy ideał $I\subseteq\Z$ ma postać $n\Z$ dla pewnego
$n\ge0$. Gdy $I=\{0\}$, bierzemy $n=0$.
W przeciwnym przypadku $I$ zawiera dodatnią liczbę całkowitą
(jeżeli $a\in I$ jest ujemne, to $-a\in I$). Niech $n$ będzie
najmniejszą dodatnią liczbą w $I$.
Oczywiście $n\Z\subseteq I$. Dla dowolnego $a\in I$ wykonujemy
dzielenie z resztą $a=qn+r$, gdzie $0\le r<n$.
Wtedy $r=a-qn\in I$. Minimalność $n$ wymusza $r=0$,
więc $a\in n\Z$. Otrzymaliśmy $I=n\Z$.

Dla $n\ge2$ ideał $(n)$ jest pierwszy (równoważnie: maksymalny)
dokładnie wtedy, gdy $n$ jest liczbą pierwszą.
Istotnie, dla liczby pierwszej $p$ iloraz $\Z/(p)$ jest ciałem,
a dla liczby złożonej $n$ ma dzielniki zera.
Ideał $(0)$ jest pierwszy, ponieważ $\Z$ jest dziedziną,
ale nie jest maksymalny, ponieważ $\Z$ nie jest ciałem.
\end{przyklad}

\begin{przyklad}[Różnica między jednym a dwiema zmiennymi]
Niech $k$ będzie ciałem. W $k[x,y]$ ideał $(x)$ jest pierwszy,
ponieważ po podstawieniu $x=0$ uzyskujemy
$k[x,y]/(x)\cong k[y]$, a $k[y]$ jest dziedziną.
Nie jest jednak maksymalny, ponieważ $k[y]$ nie jest ciałem:
wielomian $y$ nie ma odwrotności w $k[y]$.
Z kolei $(x,y)$ jest maksymalny, ponieważ
$k[x,y]/(x,y)\cong k$ (w ilorazie $x$ i $y$ stają się zerem).
Dla pierwszego izomorfizmu wystarczy rozważyć homomorfizm
$\psi\colon k[x,y]\to k[y]$ dany przez $\psi(f)=f(0,y)$.
Jest on surjektywny (każdy wielomian w $y$ jest obrazem samego siebie).
Jeśli $f(x,y)=\sum_{j=0}^m a_j(y)x^j$, to $\psi(f)=a_0(y)$,
zatem $\psi(f)=0$ wtedy i tylko wtedy, gdy $a_0(y)=0$,
czyli gdy $f$ jest wielokrotnością $x$. Mamy $\kerf\psi=(x)$,
więc $k[x,y]/(x)\cong k[y]$. Drugi izomorfizm zostanie
uzasadniony ogólnym twierdzeniem o ideale punktu.
\end{przyklad}

\subsection{Pierwiastnik ideału i informacja, której nie widać na punktach}

\begin{definicja}[Pierwiastek ideału]
Dla ideału $I\subseteq A$ definiujemy jego \emph{pierwiastek}
$\sqrt I=\{a\in A:\text{istnieje }m\ge1\text{ takie, że }a^m\in I\}$.
Ideał $I$ nazywamy \emph{radykalnym}, gdy $\sqrt I=I$.
Pierścień $A$ jest \emph{zredukowany}, gdy jego jedynym
nilpotentnym elementem jest $0$.
\end{definicja}

\begin{twierdzenie}
\label{tw:pierwiastek}
Dla każdego ideału $I$ pierścienia przemiennego $A$
zbiór $\sqrt I$ jest ideałem zawierającym $I$.
Ponadto $A/I$ jest zredukowany wtedy i tylko wtedy, gdy $I$ jest radykalny.
\end{twierdzenie}
\begin{proof}
Jeżeli $a\in I$, to $a^1\in I$, więc $I\subseteq\sqrt I$.
Jeśli $a^m\in I$ oraz $r\in A$, to $(ra)^m=r^ma^m\in I$,
a więc $ra\in\sqrt I$. Podobnie $(-a)^m=(-1)^ma^m\in I$,
więc $-a\in\sqrt I$.

Niech teraz $a,b\in\sqrt I$, czyli $a^m\in I$ i $b^n\in I$
dla pewnych $m,n\ge1$. Rozwińmy $(a+b)^{m+n-1}$
wzorem dwumianowym. W każdym składniku
$\binom{m+n-1}{j}a^jb^{m+n-1-j}$
mamy $j\ge m$ lub $m+n-1-j\ge n$:
w przeciwnym razie $j\le m-1$ i $m+n-1-j\le n-1$,
co po dodaniu dałoby $m+n-1\le m+n-2$, sprzeczność.
Każdy składnik zawiera zatem czynnik $a^m$ lub $b^n$
i należy do $I$. Suma także należy do $I$.
Stąd $a+b\in\sqrt I$, co kończy dowód, że $\sqrt I$ jest ideałem.

Element $a+I$ jest nilpotentny w $A/I$ dokładnie wtedy, gdy
$(a+I)^m=0+I$ dla pewnego $m\ge1$, a to znaczy $a^m\in I$,
czyli $a\in\sqrt I$. W ilorazie nie ma niezerowych nilpotentów
dokładnie wtedy, gdy każdy taki $a$ należy już do $I$.
Warunek ten jest równoważny $\sqrt I=I$.
\end{proof}

\begin{przyklad}[Podwójny punkt]
Niech $k$ będzie ciałem, $A=k[t]$, $I=(t^2)$.
Pierścień $A/I$ zawiera niezerowy element $\overline t=t+(t^2)$,
ale $\overline t^2=0$. Element jest niezerowy,
bo wielomian $t$ nie może być wielokrotnością $t^2$
(jego stopień jest mniejszy).
Mamy $\sqrt{(t^2)}=(t)$: jeśli $f(t)^m$ jest podzielny przez $t^2$,
to $f(0)^m=0$, więc $f(0)=0$ i $t\mid f$;
odwrotnie, dla $f=tg$ mamy $f^2=t^2g^2\in(t^2)$.

Zatem $k[t]/(t)$ i $k[t]/(t^2)$ opisują ten sam warunek
na \emph{zwykłych punktach nad ciałem}: $t=0$.
Nie są jednak izomorficzne jako pierścienie, ponieważ w drugim
występuje niezerowy nilpotent, a w pierwszym (izomorficznym z $k$) nie.
\end{przyklad}

\section{Wielomiany dokładniej: stopień, dzielenie i pierwiastki}

\begin{definicja}[Stopień wielomianu]
Niech $A$ będzie pierścieniem. Jeśli
$f=a_0+a_1x+\dots+a_dx^d\in A[x]$ i $a_d\ne0$,
to $\deg f=d$, a $a_d$ jest \emph{współczynnikiem wiodącym}.
Dla wielomianu zerowego przyjmujemy formalnie $\deg0=-\infty$.
Wielomian jest \emph{moniczny}, gdy jego współczynnik wiodący jest równy $1$.
\end{definicja}

\begin{stwierdzenie}[Stopień iloczynu]
Jeżeli $A$ jest dziedziną całkowitości, a $f,g\in A[x]$ są niezerowe,
to $\deg(fg)=\deg f+\deg g$. W szczególności $A[x]$ jest dziedziną.
\end{stwierdzenie}
\begin{proof}
Niech $f=a_mx^m+\dots$ i $g=b_nx^n+\dots$, gdzie $a_m,b_n\ne0$.
Współczynnik przy $x^{m+n}$ w iloczynie wynosi $a_mb_n$:
aby uzyskać wykładnik $m+n$, trzeba pomnożyć wyrazy najwyższych stopni.
W dziedzinie $a_mb_n\ne0$. Nie ma wyrazów wyższego stopnia,
więc $\deg(fg)=m+n$. W szczególności iloczyn dwóch niezerowych
wielomianów nie jest zerowy, a zatem $A[x]$ jest dziedziną.
\end{proof}

\begin{uwaga}
Założenia o dziedzinie nie wolno pomijać. W $(\Z/6\Z)[x]$
mamy $(\overline2x)(\overline3x)=0$, choć oba czynniki są
niezerowymi wielomianami stopnia $1$.
\end{uwaga}

\begin{twierdzenie}[Dzielenie z resztą]
\label{tw:dzielenie}
Niech $k$ będzie ciałem, a $f,g\in k[x]$, przy czym $g\ne0$.
Istnieją dokładnie jedne wielomiany $q,r\in k[x]$ takie, że
$f=qg+r$ oraz $r=0$ lub $\deg r<\deg g$.
\end{twierdzenie}
\begin{proof}
Najpierw konstruujemy $q$ i $r$.
Jeśli $f=0$ albo $\deg f<\deg g$, przyjmujemy $q=0,r=f$.
W przeciwnym przypadku zapiszmy
$f=ax^m+\text{wyrazy niższe}$,
$g=bx^n+\text{wyrazy niższe}$, gdzie $m\ge n$ i $a,b\ne0$.
Ponieważ $k$ jest ciałem, element $b$ ma odwrotność.
Wielomian $h=f-ab^{-1}x^{m-n}g$ ma stopień mniejszy od $m$:
wyrazy wiodące obu składników dokładnie się znoszą.
Powtarzając tę operację na $h$, po skończonej liczbie kroków
(za każdym razem stopień maleje) otrzymujemy resztę $r$ o stopniu
mniejszym od $\deg g$ lub zero, a zebrane wielomiany
$ab^{-1}x^{m-n}$ tworzą iloraz $q$. To dowodzi istnienia.

Dla jedyności załóżmy, że $f=qg+r=q'g+r'$ i obie reszty
mają stopień mniejszy od $\deg g$ albo są zerowe.
Mamy $(q-q')g=r'-r$. Jeśli $q-q'\ne0$, to z własności stopnia
w dziedzinie $k[x]$ dostajemy
$\deg((q-q')g)=\deg(q-q')+\deg g\ge\deg g$.
Z drugiej strony $\deg(r'-r)<\deg g$ albo $r'-r=0$,
co jest niemożliwe. Zatem $q=q'$, a w konsekwencji $r=r'$.
\end{proof}

\begin{przyklad}[Dzielenie w praktyce]
W $\Q[x]$ podzielmy $f=x^3-2x+3$ przez $g=x-2$.
Najpierw odejmujemy $x^2(x-2)=x^3-2x^2$,
co zostawia $2x^2-2x+3$.
Następnie odejmujemy $2x(x-2)=2x^2-4x$,
co zostawia $2x+3$. Odejmujemy $2(x-2)=2x-4$
i otrzymujemy resztę $7$. A zatem
$x^3-2x+3=(x^2+2x+2)(x-2)+7$.
Warto sprawdzić wynik przez podstawienie $x=2$:
$f(2)=8-4+3=7$, zgodnie z resztą.
\end{przyklad}

\begin{wniosek}[Twierdzenie o reszcie i czynniku]
Dla ciała $k$, wielomianu $f\in k[x]$ i elementu $a\in k$
reszta z dzielenia $f$ przez $x-a$ jest stałą $f(a)$.
W szczególności $f(a)=0$ wtedy i tylko wtedy, gdy $(x-a)\mid f$.
\end{wniosek}
\begin{proof}
Z twierdzenia o dzieleniu mamy $f=(x-a)q+r$,
gdzie $r$ jest stałą lub zerem. Podstawiając $x=a$,
dostajemy $f(a)=r$. Stąd od razu wynika równoważność.
\end{proof}

\begin{twierdzenie}[Ograniczenie liczby pierwiastków]
\label{tw:korzenie}
Jeżeli $k$ jest ciałem, a $f\in k[x]$ jest niezerowym wielomianem
stopnia $d$, to $f$ ma w $k$ co najwyżej $d$ różnych pierwiastków.
\end{twierdzenie}
\begin{proof}
Dowodzimy indukcyjnie względem $d$.
Wielomian stopnia $0$ jest niezerową stałą i nie ma pierwiastków.
Jeśli $f$ stopnia $d\ge1$ nie ma pierwiastków, teza jest oczywista.
W przeciwnym razie wybierzmy pierwiastek $a$.
Na mocy twierdzenia o czynniku $f=(x-a)g$, gdzie $\deg g=d-1$.
Dla każdego innego pierwiastka $b\ne a$ zachodzi
$0=f(b)=(b-a)g(b)$.
Ponieważ $k$ jest ciałem, $b-a\ne0$ i możemy skrócić ten czynnik,
więc $g(b)=0$. Pozostałe pierwiastki $f$ są więc pierwiastkami $g$,
których jest najwyżej $d-1$ z założenia indukcyjnego.
Łącznie $f$ ma najwyżej $d$ różnych pierwiastków.
\end{proof}

\begin{uwaga}
Nie jest to prawda dla wielomianów nad dowolnymi pierścieniami:
w $\Z/8\Z$ wielomian $x^2-1$ ma cztery różne pierwiastki,
co widzieliśmy na początku. To dlatego warunek „$k$ jest ciałem”
nie jest kosmetyczny.
\end{uwaga}

\subsection{Algorytm Euklidesa i ideały w \texorpdfstring{$k[x]$}{k[x]}}

\begin{definicja}[Podzielność i element nierozkładalny]
Niech $A$ będzie dziedziną. Dla $a,b\in A$ piszemy $a\mid b$,
gdy istnieje $c\in A$ takie, że $b=ac$.
Element $u\in A$ jest \emph{nierozkładalny}, jeśli jest niezerowy,
nie jest jednostką i z $u=ab$ wynika, że $a$ lub $b$ jest jednostką.
Dwa elementy $a,b$ nazywamy \emph{stowarzyszonymi}, jeśli
$a=ub$ dla pewnej jednostki $u\in A^\times$.
\end{definicja}

\begin{stwierdzenie}[Ideały w {$k[x]$} są główne]
\label{st:kx-glowne}
Jeżeli $k$ jest ciałem, to każdy ideał $I\subseteq k[x]$
ma postać $(d)$ dla pewnego $d\in k[x]$.
Dla niezerowego $I$ wielomian $d$ można wybrać moniczny,
a wtedy jest on jednoznacznie wyznaczony.
\end{stwierdzenie}
\begin{proof}
Dla $I=(0)$ bierzemy $d=0$.
Jeżeli $I\ne(0)$, wybieramy niezerowy wielomian $d\in I$
o najmniejszym możliwym stopniu. Mnożąc przez odwrotność
współczynnika wiodącego, możemy założyć, że $d$ jest moniczny.
Dla dowolnego $f\in I$ dzielimy przez $d$:
$f=qd+r$, gdzie $r=0$ lub $\deg r<\deg d$.
Ponieważ $f,qd\in I$, mamy $r=f-qd\in I$.
Minimalność stopnia $d$ wymusza $r=0$, czyli $d\mid f$.
Stąd $I\subseteq(d)$. Druga inkluzja wynika z $d\in I$.

Jeśli moniczne $d,e$ generują ten sam niezerowy ideał,
to $d\mid e$ i $e\mid d$. Z wzoru na stopień otrzymujemy
$\deg d=\deg e$, więc $e=cd$ dla stałej $c\in k^\times$.
Moniczność obu wielomianów wymusza $c=1$, czyli $d=e$.
\end{proof}

\begin{twierdzenie}[Wielomianowy Bézout]
\label{tw:bezout-pol}
Niech $k$ będzie ciałem, a $f,g\in k[x]$ nie będą równocześnie zerowe.
Istnieje jedyny moniczny wielomian $d\in k[x]$ taki, że
$(f,g)=(d)$. Ponadto istnieją $u,v\in k[x]$ takie, że
$d=uf+vg$. Wielomian $d$ jest największym wspólnym dzielnikiem
$f$ i $g$ w sensie podzielności, znormalizowanym do postaci monicznej.
\end{twierdzenie}
\begin{proof}
Z twierdzenia~\ref{st:kx-glowne} ideał $(f,g)$ ma jedyny moniczny
generator $d$. Ponieważ $f,g\in(d)$, mamy $d\mid f$ oraz $d\mid g$.
Każdy wspólny dzielnik $c$ wielomianów $f,g$ dzieli wszystkie
kombinacje $uf+vg$; w szczególności dzieli $d$.
Ponieważ $d\in(f,g)$, na mocy jawnego opisu generowania
istnieją $u,v\in k[x]$ takie, że $d=uf+vg$.
\end{proof}

\begin{przyklad}[Obliczenie NWD wielomianów]
W $\Q[x]$ rozważmy $f=x^3-1$ i $g=x^2-1$.
Mamy $f=xg+(x-1)$ oraz $g=(x+1)(x-1)$,
więc $\gcd(f,g)=x-1$ (w postaci monicznej).
Co więcej, bezpośrednio z pierwszego dzielenia dostajemy
tożsamość Bézouta $x-1=f-xg$.
To świadectwo mówi, że ideał $(f,g)$ jest równy $(x-1)$.
\end{przyklad}

\begin{twierdzenie}[Ideały maksymalne w {$k[x]$}]
\label{tw:irreducible}
Niech $k$ będzie ciałem, a $f\in k[x]$ wielomianem niezerowym
i niestałym. Ideał $(f)$ jest maksymalny wtedy i tylko wtedy,
gdy $f$ jest nierozkładalny w $k[x]$. W takim przypadku
$k[x]/(f)$ jest ciałem.
\end{twierdzenie}
\begin{proof}
Załóżmy, że $f$ jest nierozkładalny.
Niech $g+(f)$ będzie niezerowym elementem ilorazu.
Wtedy $f\nmid g$. Każdy wspólny dzielnik $f$ i $g$
jest albo jednostką, albo jest stowarzyszony z $f$,
gdyż $f$ jest nierozkładalny; druga możliwość dawałaby $f\mid g$.
Zatem $\gcd(f,g)=1$. Z twierdzenia~\ref{tw:bezout-pol}
istnieją $u,v\in k[x]$ takie, że $uf+vg=1$.
Po przejściu do ilorazu dostajemy
$(v+(f))(g+(f))=1+(f)$. Każda niezerowa klasa ma więc odwrotność,
zatem $k[x]/(f)$ jest ciałem, a $(f)$ maksymalny.

Odwrotnie, jeśli $f=gh$ dla niejednostek $g,h\in k[x]$,
to $\deg g,\deg h\ge1$ i oba stopnie są mniejsze od $\deg f$.
Mamy $(f)\subseteq(g)$. Inkluzja jest ścisła:
gdyby $g\in(f)$, to $g=qf$ dla niezerowego $q$,
a więc $\deg g\ge\deg f$, sprzeczność.
Ideał $(g)$ jest właściwy, bo $g$ nie jest jednostką.
W takim razie $(f)$ nie jest maksymalny.
\end{proof}

\begin{przyklad}[Konstrukcja ciała liczb zespolonych]
Wielomian $x^2+1\in\R[x]$ jest nierozkładalny nad $\R$:
gdyby rozkładał się na dwa wielomiany liniowe,
miałby rzeczywisty pierwiastek, a $a^2+1>0$ dla każdego $a\in\R$.
Dlatego $\R[x]/(x^2+1)$ jest ciałem.
Jeśli $i=x+(x^2+1)$, to $i^2=-1$, a każdą klasę można zapisać
jednoznacznie jako $a+bi$ dla $a,b\in\R$,
ponieważ dzielenie przez $x^2+1$ zostawia resztę stopnia co najwyżej $1$.
Działania zgadzają się z działaniami w $\C$,
więc $\R[x]/(x^2+1)\cong\C$.
\end{przyklad}

\section{Wiele zmiennych i pierwsze pierścienie współrzędnych}

\subsection{Podstawianie, ideały punktów i usuwanie zmiennych}

Niech $k$ będzie ciałem. Pierścień $k[x_1,\dots,x_n]$
można skonstruować iteracyjnie:
$k[x_1,\dots,x_n]=k[x_1,\dots,x_{n-1}][x_n]$.
Elementami są skończone sumy jednomianów
$c_{\alpha}x_1^{\alpha_1}\cdots x_n^{\alpha_n}$,
gdzie $\alpha\in\Z_{\ge0}^n$ i $c_\alpha\in k$.
Równość wielomianów oznacza równość współczynników.

\begin{stwierdzenie}[Podstawienie do punktu]
\label{st:eval-point}
Dla $a=(a_1,\dots,a_n)\in k^n$ homomorfizm
$\ev_a\colon k[x_1,\dots,x_n]\to k$ zadany przez
$\ev_a(f)=f(a_1,\dots,a_n)$ jest surjektywny, a jego jądro wynosi
$\mathfrak m_a=(x_1-a_1,\dots,x_n-a_n)$.
W szczególności $\mathfrak m_a$ jest ideałem maksymalnym.
\end{stwierdzenie}
\begin{proof}
Homomorfizm jest surjektywny, bo każda stała $c\in k$
jest obrazem stałego wielomianu $c$. Każdy generator $x_j-a_j$
przechodzi na zero, więc
$\mathfrak m_a\subseteq\kerf\ev_a$.

Dla odwrotnej inkluzji weźmy wielomian $f$ spełniający $f(a)=0$.
Kolejno zastępujmy zmienne ich wartościami:
\begin{align*}
 f(x_1,\dots,x_n)-f(a_1,\dots,a_n)
 &=\bigl[f(x_1,\dots,x_n)-f(a_1,x_2,\dots,x_n)\bigr]\\
 &\quad+\bigl[f(a_1,x_2,\dots,x_n)-f(a_1,a_2,x_3,\dots,x_n)\bigr]
 +\dots\\
 &\quad+\bigl[f(a_1,\dots,a_{n-1},x_n)-f(a_1,\dots,a_n)\bigr].
\end{align*}
Każdy nawias jest podzielny przez odpowiedni $x_j-a_j$:
to wzór $x_j^m-a_j^m=(x_j-a_j)\sum_{r=0}^{m-1}x_j^{m-1-r}a_j^r$
zastosowany współczynnik po współczynniku w pozostałych zmiennych.
Zatem $f-f(a)\in\mathfrak m_a$. Skoro $f(a)=0$,
mamy $f\in\mathfrak m_a$. Otrzymaliśmy równość jądra i ideału.
Ponieważ iloraz przez jądro jest izomorficzny z ciałem $k$,
ideał $\mathfrak m_a$ jest maksymalny na mocy
 twierdzenia~\ref{tw:kryterium}.
\end{proof}

\begin{uwaga}[Co wynika, a co jeszcze nie wynika?]
Udowodniliśmy, że każdemu punktowi $a\in k^n$ odpowiada ideał
maksymalny $\mathfrak m_a$. Nie udowodniliśmy odwrotnego zdania,
że każdy ideał maksymalny ma tę postać. Jest ono fałszywe dla
niektórych ciał, np. nad $\R$ ideał $(x^2+1)\subseteq\R[x]$
jest maksymalny, lecz nie pochodzi od żadnego punktu $\R$.
Dla ciała algebraicznie domkniętego odwrotne zdanie jest częścią
słabego twierdzenia Hilberta o zerach, którym zajmiemy się
w rozdziale poświęconym zbiorom algebraicznym.
\end{uwaga}

\begin{przyklad}[Parabola i jej pierścień]
\label{pr:parabola}
Rozważmy homomorfizm
$\Phi\colon k[x,y]\to k[t]$ zadany przez
$\Phi(f(x,y))=f(t,t^2)$.
Jest on surjektywny: $t=\Phi(x)$, więc każdy wielomian w $t$
jest obrazem odpowiedniego wielomianu w $x$.
Ponadto $\Phi(y-x^2)=t^2-t^2=0$, czyli
$(y-x^2)\subseteq\kerf\Phi$.

Aby udowodnić inkluzję odwrotną, dla dowolnego
$f(x,y)=\sum_{j=0}^m a_j(x)y^j$ zapiszmy
\[
 f(x,y)-f(x,x^2)
 =\sum_{j=1}^m a_j(x)(y^j-x^{2j})
 =(y-x^2)\sum_{j=1}^m a_j(x)
   \sum_{r=0}^{j-1}y^{j-1-r}x^{2r}.
\]
Jeżeli $\Phi(f)=0$, to $f(t,t^2)=0$ jako wielomian w $k[t]$,
a więc również $f(x,x^2)=0$ jako wielomian w $k[x]$
(oznaczenie zmiennej nie ma znaczenia). Z powyższego wzoru
wynika $f\in(y-x^2)$. Zatem $\kerf\Phi=(y-x^2)$.
Pierwsze twierdzenie o izomorfizmie daje
\[
 k[x,y]/(y-x^2)\cong k[t].
\]
Algebraicznie oznacza to, że po narzuceniu relacji $y=x^2$
współrzędna $y$ jest już określona przez $x$: każda klasa wielomianów
ma reprezentanta zależnego tylko od $x$.
Nad nieskończonym ciałem $k$ izomorfizm można też interpretować
bezpośrednio jako opis funkcji wielomianowych na punktach paraboli.
Nad ciałem skończonym trzeba zachować rozróżnienie z
uwagi~\ref{uw:formalna-funkcja}.
\end{przyklad}

\begin{center}
\begin{tikzpicture}[scale=0.86,>=Stealth]
\draw[->,gray!80] (-2.6,0)--(2.7,0) node[right] {$x$};
\draw[->,gray!80] (0,-0.25)--(0,4.9) node[above] {$y$};
\draw[blue!70!black,very thick,domain=-2.12:2.12,samples=100]
  plot (\x,{\x*\x});
\fill[red!75!black] (1.35,1.8225) circle (2pt);
\draw[densely dashed,red!75!black] (1.35,0)--(1.35,1.8225)--(0,1.8225);
\node[below,red!75!black] at (1.35,0) {$t$};
\node[left,red!75!black] at (0,1.8225) {$t^2$};
\node[blue!70!black,anchor=west] at (1.25,4.1) {$y=x^2$};
\end{tikzpicture}
\end{center}
Punkt paraboli opisujemy jedną liczbą $t$ poprzez parametryzację
$t\mapsto(t,t^2)$. Obliczenie jądra $\Phi$ jest algebraicznym
odpowiednikiem tego rysunku, ale działa nad \emph{dowolnym} ciałem $k$,
również wtedy, gdy nie możemy rysować jego punktów na płaszczyźnie.

\begin{przyklad}[Dwie osie i iloraz przez $(xy)$]
W $k[x,y]$ rozważmy ideał $(xy)$.
Równanie $xy=0$ opisuje sumę dwóch osi: $x=0$ lub $y=0$.
Pierścień $k[x,y]/(xy)$ pamięta, że iloczyn współrzędnych $x$ i $y$
ma być zerem, choć żadna z tych współrzędnych nie musi być zerem
jako element tego pierścienia.
W istocie ich klasy są niezerowe. Gdyby $x=xyh$ dla pewnego
$h\in k[x,y]$, to po skróceniu niezerowego $x$ w dziedzinie
$k[x,y]$ otrzymalibyśmy $1=yh$, co po podstawieniu $y=0$
dawałoby sprzeczność $1=0$ w ciele $k$. Zatem $x\notin(xy)$;
analogicznie $y\notin(xy)$. Jednocześnie
$(x+(xy))(y+(xy))=xy+(xy)=0+(xy)$.
Iloraz ma zatem dzielniki zera. Jest to algebraiczny ślad faktu,
że zbiór geometryczny składa się z dwóch składowych.
\end{przyklad}

\begin{center}
\begin{tikzpicture}[scale=0.87,>=Stealth]
\draw[gray!65,->] (-2.4,0)--(2.5,0) node[right] {$x$};
\draw[gray!65,->] (0,-2.2)--(0,2.5) node[above] {$y$};
\draw[blue!70!black,line width=2.2pt] (-2.1,0)--(2.1,0);
\draw[red!75!black,line width=2.2pt] (0,-1.95)--(0,2.05);
\fill[black] (0,0) circle (2.3pt);
\node[blue!70!black,below right] at (1.25,0) {$y=0$};
\node[red!75!black,above right] at (0,1.2) {$x=0$};
\node[align=left,font=\small] at (3.9,-1.3)
  {Przekrój rzeczywisty\\zbioru $V(xy)$.};
\end{tikzpicture}
\end{center}

\begin{przyklad}[Pierścień wielomianów w dwóch zmiennych nie jest pierścieniem ideałów głównych]
Przypuśćmy, że w $k[x,y]$ ideał $(x,y)$ jest równy $(f)$
dla pewnego $f\in k[x,y]$. Wtedy $f\mid x$ i $f\mid y$.
Pokażemy, że każdy wspólny dzielnik $x$ i $y$ jest jednostką.
Jeżeli $x=fg$ w $k[x,y]$, traktujemy wszystkie wielomiany
jako wielomiany w $x$ o współczynnikach w dziedzinie $k[y]$.
Stopnie względem $x$ sumują się, więc jeden z $f,g$ ma stopień
względem $x$ równy zero. Jeżeli $f\in k[y]$, to z $x=fg$
współczynnik przy $x$ pokazuje, że $f$ dzieli $1$ w $k[y]$,
więc $f$ jest jednostką. W przeciwnym razie $g\in k[y]$
i ten sam argument pokazuje, że $g$ jest jednostką;
wtedy $f$ jest stowarzyszony z $x$.
W rezultacie każdy dzielnik $x$ jest jednostką albo jest
stowarzyszony z $x$.

Gdyby $f$ nie był jednostką, musiałby być stowarzyszony z $x$,
ale z $f\mid y$ wynikałoby wtedy $x\mid y$, co jest niemożliwe:
po podstawieniu $x=0$ równanie $y=xh(x,y)$ dawałoby $y=0$
w $k[y]$. Zatem $f$ jest jednostką, skąd $(f)=k[x,y]$.
Tymczasem $(x,y)$ jest ideałem właściwym, ponieważ
wszystkie jego elementy przyjmują wartość zero w $(0,0)$,
a wielomian $1$ nie. Sprzeczność dowodzi, że $(x,y)$ nie jest główny.
\end{przyklad}

\section{Czego uczą nas przykłady geometryczne?}

Przyjrzyjmy się trzem ilorazom pierścienia $k[x,y]$:
\begin{enumerate}[label=\arabic*)]
\item $k[x,y]/(y-x^2)$ odpowiada paraboli. Jest izomorficzny z $k[t]$,
więc jest dziedziną: ideał $(y-x^2)$ jest pierwszy.
Gdy $k$ jest algebraicznie domknięte, oznacza to też
nierozkładalność paraboli w topologii Zariskiego; dokładne uzasadnienie
podamy w następnym rozdziale. Nie wyciągamy tego wniosku bez założeń
o ciele: nad ciałem skończonym topologia Zariskiego na zbiorze
punktów $k^n$ jest dyskretna.
\item $k[x,y]/(xy)$ odpowiada sumie osi. Ma dzielniki zera,
bo klasy $x$ i $y$ są niezerowe, ale ich iloczyn jest zerem.
\item $k[t]/(t^2)$ ma niezerowy nilpotent. W odróżnieniu od
$k[t]/(t)$ przechowuje informację o „pogrubieniu” punktu $t=0$.
\end{enumerate}
Te obserwacje prowadzą do trzech idei, które rozwinie kolejny rozdział.

\textbf{Pierwsza idea:} badamy nie tylko rozwiązania $f=0$, ale cały
ideał równań, ponieważ kombinacje równań również muszą znikać.

\textbf{Druga idea:} iloraz pierścienia wielomianów przez ideał równań
koduje algebraiczne funkcje na rozpatrywanym obiekcie.

\textbf{Trzecia idea:} własności ideału odzwierciedlają geometrię.
Ideał pierwszy daje iloraz bez dzielników zera,
a pierwiastek ideału kontroluje, czy iloraz zawiera nilpotenty.
Dopiero nad ciałem algebraicznie domkniętym połączymy pierwszość
ideału z nierozkładalnością geometryczną odpowiadającego mu zbioru.

\begin{uwaga}[Granica dotychczasowej teorii]
W tym rozdziale nie dowodziliśmy jeszcze, że dla dowolnego ideału
$I\subseteq k[x_1,\dots,x_n]$ odpowiednie punkty geometryczne
pozwalają odzyskać $\sqrt I$. Takie twierdzenie wymaga
założenia, że $k$ jest algebraicznie domknięte, i stanowi treść
silnego twierdzenia Hilberta o zerach. Będzie to główny cel
pierwszego rozdziału właściwej geometrii algebraicznej.
\end{uwaga}

\section{Podsumowanie logicznych zależności}

\begin{center}
\begin{tikzpicture}[>=Stealth,
  box/.style={draw=gray!65,rounded corners=3pt,align=center,inner sep=6pt,fill=gray!5},
  node distance=9mm and 14mm]
\node[box] (r) {pierścień $A$};
\node[box,below left=of r] (i) {ideał $I\subseteq A$};
\node[box,below right=of r] (h) {homomorfizm\\$\varphi\colon A\to B$};
\node[box,below=of i] (q) {iloraz $A/I$};
\node[box,below=of h] (k) {jądro $\kerf\varphi$};
\node[box,below=of q] (p) {ideał pierwszy\\$\Longleftrightarrow A/I$ dziedzina};
\node[box,below=of k] (m) {ideał maksymalny\\$\Longleftrightarrow A/I$ ciało};
\draw[->] (r)--(i);
\draw[->] (r)--(h);
\draw[->] (i)--(q);
\draw[->] (h)--(k);
\draw[->] (k)--(q);
\draw[->] (q)--(p);
\draw[->] (q)--(m);
\end{tikzpicture}
\end{center}

Ten diagram nie oznacza, że każdy ideał jest pierwszy lub maksymalny.
Dwie dolne ramki podają \emph{kryteria} dla ideałów, które mają
daną własność. Diagram przypomina natomiast przepływ pojęć:
ideał pozwala utworzyć iloraz, a typ ilorazu rozstrzyga o typie ideału.

W dalszej części skryptu przejdziemy do zbiorów algebraicznych
w $k^n$, topologii Zariskiego, korespondencji między zbiorami
a ideałami oraz twierdzenia Hilberta o zerach. Wyjaśnimy też,
czym różni się zbiór punktów nad $k$ od jego pierścienia współrzędnych. Następnie
omówimy geometrię rzutową, osobliwości, normalizację i krzywe
algebraiczne wraz z rodzajem oraz twierdzeniem Riemanna--Rocha.

\section{Jednoznaczność rozkładu wielomianów na czynniki}
\label{sec:ufd-preparation}

W dalszej geometrii będziemy używać informacji silniejszej niż
algorytm Euklidesa dla $k[t]$: wielomiany w \emph{wielu} zmiennych
również mają jednoznaczny rozkład na czynniki nierozkładalne.
Trzeba to uzasadnić; nie wynika to automatycznie z faktu,
że $k[x,y]$ jest dziedziną.

\begin{definicja}[Pierścień z jednoznacznym rozkładem]
Niech $A$ będzie dziedziną całkowitości. Elementy $a,b\in A$
są \emph{stowarzyszone}, gdy $a=ub$ dla pewnej jednostki
$u\in A^\times$. Dziedzinę $A$ nazywamy \emph{pierścieniem
z jednoznacznym rozkładem na czynniki} (ang. \emph{UFD}), jeżeli
każdy niezerowy nieodwracalny element jest skończonym iloczynem
elementów nierozkładalnych, a dwa takie rozkłady różnią się
tylko kolejnością czynników i ich stowarzyszeniem.
\end{definicja}

\begin{lemat}[Lemat Gaussa]
\label{lem:gauss-ufd}
Niech $A$ będzie pierścieniem z jednoznacznym rozkładem,
$K=\Frac(A)$, a $f,g\in A[T]$. Wielomian jest
\emph{pierwotny}, gdy żaden nieodwracalny element $A$
nie dzieli wszystkich jego współczynników.
Wówczas: (i) iloczyn dwóch pierwotnych wielomianów jest
pierwotny; (ii) jeśli pierwotny $f$ rozkłada się
w $K[T]$ na czynniki dodatnich stopni, to rozkłada się
także w $A[T]$ na czynniki dodatnich stopni.
\end{lemat}
\begin{proof}
(i) Gdyby iloczyn $fg$ nie był pierwotny, to wszystkie
jego współczynniki dzieliłby pewien element nierozkładalny
$p\in A$. W pierścieniu z jednoznacznym rozkładem każdy
element nierozkładalny jest pierwszy: jeśli $p\mid ab$,
porównanie rozkładów na czynniki wymusza $p\mid a$ albo
$p\mid b$. Zatem $A/(p)$ jest dziedziną.
Po redukcji współczynników modulo $p$ niezerowe
wielomiany $\bar f,\bar g\in(A/(p))[T]$ miałyby
iloczyn zero, co przeczy temu, że pierścień
wielomianów nad dziedziną jest dziedziną.

(ii) Załóżmy $f=uv$ dla $u,v\in K[T]$ o dodatnich stopniach.
Usuwając mianowniki i wyciągając wspólne czynniki
współczynników, możemy przedstawić
$u=\alpha u_0$, $v=\beta v_0$, gdzie
$\alpha,\beta\in K^\times$, a $u_0,v_0\in A[T]$
są pierwotne. Z (i) iloczyn $u_0v_0$ jest pierwotny.
Mamy więc $f=(\alpha\beta)u_0v_0$.
Jeśli iloczyn dwóch pierwotnych wielomianów jest mnożony
przez ułamek $a/b\in K$ w postaci nieskracalnej,
a wynik ma współczynniki w $A$, to $b$ musi być jednostką:
w przeciwnym razie pewien czynnik nierozkładalny $p\mid b$
dzieliłby wszystkie współczynniki $a(u_0v_0)$;
przez $p\nmid a$ dzieliłby wszystkie współczynniki
$u_0v_0$, sprzeczność. Po wykazaniu, że $b$ jest jednostką, skalar $a/b$ należy do $A$.
Gdyby nie był jednostką, miałby w $A$ dzielnik nierozkładalny
$p$, który dzieliłby wszystkie współczynniki $f=(a/b)u_0v_0$.
Przeczyłoby to pierwotności $f$. Zatem $a/b\in A^\times$.
W konsekwencji $f$ jest iloczynem dwóch wielomianów
w $A[T]$ dodatnich stopni.
\end{proof}

\begin{twierdzenie}[Pierścienie wielomianów są UFD]
\label{tw:polynomial-UFD}
Jeśli $k$ jest ciałem i $n\ge1$, to
$k[x_1,\dots,x_n]$ jest pierścieniem z jednoznacznym
rozkładem na czynniki. Bardziej ogólnie, jeśli $A$
jest UFD, to $A[T]$ jest UFD.
\end{twierdzenie}
\begin{proof}
Najpierw udowodnimy zdanie ogólne. Dla niezerowego
$f\in A[T]$ wyznaczmy \emph{zawartość} $c(f)$, czyli
największy wspólny dzielnik jego niezerowych współczynników,
określony z dokładnością do jednostki. Dzięki istnieniu
i jednoznaczności rozkładu współczynników w $A$ taki
element istnieje. Mamy $f=c(f)f_0$ z $f_0$ pierwotnym.
Lemat Gaussa pokazuje, że iloczyn dwóch pierwotnych
wielomianów jest pierwotny, a zatem
$c(fg)$ jest stowarzyszony z $c(f)c(g)$.

W ciele ułamków $K$ pierścień $K[T]$ ma jednoznaczny rozkład
na czynniki nierozkładalne. Uzasadnijmy to na podstawie
wcześniejszych wyników. Algorytm dzielenia z resztą
(twierdzenie~\ref{tw:dzielenie}) pokazuje, że każdy ideał
$K[T]$ jest główny (stwierdzenie~\ref{st:kx-glowne}).
Jeśli nierozkładalny $p$ dzieli $ab$, lecz nie dzieli $a$,
to największy wspólny dzielnik $p,a$ jest równy $1$
z dokładnością do jednostki. Tożsamość Bézouta
(twierdzenie~\ref{tw:bezout-pol}) daje $up+va=1$.
Po pomnożeniu przez $b$ otrzymujemy $upb+vab=b$;
oba składniki po lewej są podzielne przez $p$, zatem $p\mid b$.
Każdy nierozkładalny jest więc pierwszy. Rozkład niezerowego
wielomianu niebędącego jednostką istnieje przez indukcję względem
stopnia: wielomian nierozkładalny kończy postępowanie, a rozkładalny
jest iloczynem dwóch wielomianów dodatnich, ściśle mniejszych stopni.
Jednoznaczność wynika z pierwszości czynników: pierwszy czynnik
jednego rozkładu dzieli któryś czynnik drugiego, jest z nim
stowarzyszony, po czym można go skrócić i powtórzyć rozumowanie.
Jeżeli $f_0$ jest stały, to z pierwotności wynika,
że jest jednostką $A$, więc nie wymaga dalszego rozkładu.
W przeciwnym razie niech $f_0=q_1\cdots q_r$ będzie
rozkładem pierwotnego wielomianu $f_0$ na nierozkładalne
czynniki w $K[T]$
(z ewentualnym czynnikiem stałym włączonym do $q_1$).
Dla każdego $i$ usuwamy mianowniki współczynników $q_i$,
a następnie dzielimy przez ich największy wspólny dzielnik.
Otrzymujemy $q_i=\alpha_i h_i$, gdzie
$\alpha_i\in K^\times$, a $h_i\in A[T]$ jest pierwotny.
Z pierwszej części lematu Gaussa iloczyn $H=h_1\cdots h_r$
jest pierwotny. Ponieważ $f_0=(\prod_i\alpha_i)H$
i oba wielomiany należą do $A[T]$ oraz są pierwotne,
argument o nieskracalnym mianowniku z dowodu lematu Gaussa
pokazuje, że $\prod_i\alpha_i$ jest jednostką $A$.
Włączamy tę jednostkę do $h_1$ i uzyskujemy rozkład $f_0$
w $A[T]$. Każdy $h_i$ jest nierozkładalny w $A[T]$:
gdyby $h_i=ab$ z niejednostkami $a,b\in A[T]$, to
pierwotność $h_i$ i wzór na zawartość wymuszałyby
pierwotność obu czynników. Żaden z nich nie może być
nieodwracalną stałą, a więc oba miałyby dodatni stopień.
Dawałoby to rozkład $q_i$ w $K[T]$, sprzeczny z jego
nierozkładalnością. Istnienie rozkładu w $A[T]$ wynika
teraz z rozkładu zawartości $c(f)$ w $A$ i powyższego
rozkładu pierwotnej części $f_0$.

Dla jednoznaczności porównajmy dwa rozkłady $f$.
Po przejściu do $K[T]$ elementy stałe z $A$ stają się
jednostkami, a pierwotne czynniki nierozkładalne
pozostają nierozkładalne na mocy lematu Gaussa.
Jednoznaczność w $K[T]$ pozwala dopasować czynniki
wielomianowe parami z dokładnością do skalarów z $K$.
Uzasadnijmy dokładnie, że dwa pierwotne wielomiany
$h,h'\in A[T]$, dla których $h'=\lambda h$ z
$\lambda\in K^\times$, różnią się jednostką $A$.
Zapiszmy $\lambda=a/b$, gdzie $a,b\in A$ są względnie
pierwsze (nie mają wspólnego czynnika nierozkładalnego).
Z równości $bh'=ah$ wynika, że gdyby nierozkładalny
$p$ dzielił $b$, to przez pierwszość $p$ i $p\nmid a$
dzieliłby każdy współczynnik $h$, wbrew pierwotności.
Zatem $b$ jest jednostką, czyli $\lambda\in A$.
Gdyby $\lambda$ nie był jednostką, pewien nierozkładalny
dzielnik $\lambda$ dzieliłby wszystkie współczynniki
$h'$, znów wbrew pierwotności. Stąd
$\lambda\in A^\times$.

W każdym rozkładzie $f$ czynniki stałe nierozkładalne
pochodzą, z dokładnością do jednostek, z rozkładu
zawartości $c(f)$ w $A$, gdyż zawartość iloczynu
pierwotnych czynników jest jednostką.
Jednoznaczność rozkładu $c(f)$ w $A$ oraz wykazane
dopasowanie pierwotnych czynników wielomianowych
kończą dowód jednoznaczności w $A[T]$.

Ciało $k$ jest UFD w sposób trywialny, bo każdy jego
niezerowy element jest jednostką. Stosując wykazane
zdanie kolejno do $k[x_1]$, $k[x_1][x_2]$ itd.,
otrzymujemy tezę dla dowolnego $n$.
\end{proof}

\begin{uwaga}[Dlaczego potrzebujemy tego twierdzenia?]
W pierścieniu $k[X,Y,Z]$ będziemy korzystać z tego,
że jeśli wielomiany $F,G$ nie mają wspólnego czynnika
nierozkładalnego, a $F\mid Gh$, to $F\mid h$.
Rzeczywiście, gdy $F\ne0$, rozkładamy $F$ na czynniki
nierozkładalne. Każdy z nich jest pierwszy, ponieważ
pierścień jest UFD. Żaden nie dzieli $G$, więc z podzielności
$F\mid Gh$ wynika, że dzieli $h$ z co najmniej taką samą
krotnością, z jaką występuje w $F$: usuwamy kolejno jego
wystąpienia z iloczynu $Gh$, korzystając za każdym razem
z pierwszości i z faktu, że nie dzieli $G$.
Jednoznaczność rozkładu daje $F\mid h$.
Dla $F=0$ założenie $0\mid Gh$ oznacza $Gh=0$;
ponieważ $F$ i $G$ nie mają wspólnego czynnika
nierozkładalnego, $G\ne0$, zatem $h=0$ i także $0\mid h$.
Podobne rozumowania są podstawą dowodu twierdzenia Bézouta.
\end{uwaga}

